\documentclass[12pt]{article}

\usepackage[T2A]{fontenc}
\usepackage[utf8]{inputenc}
\DeclareUnicodeCharacter{03B6}{\ensuremath{\zeta}}
\usepackage[russian]{babel}

\usepackage{amsmath,amssymb}
\usepackage{geometry}
\geometry{margin=25mm}

\usepackage[unicode,hidelinks]{hyperref}
\pdfstringdefDisableCommands{%
  \def\mathbb#1{#1}%
  \def\alpha{alpha}%
  \def\pi{pi}%
  \def\delta{delta}%
  \def\times{x}%
  \def\zeta{zeta}%
}

\title{Вывод постоянной тонкой структуры из геометрии компактного внутреннего пространства $\mathbb{RP}^3\times S^1$}
\author{}
\date{}

\begin{document}
\maketitle

\begin{abstract}
Рассматривается гипотеза, что значение $\alpha^{-1}$ может быть связано с евклидовым эффективным действием калибровочного сектора при компактификации на $K=\mathbb{RP}^3\times S^1$.
Цель --- сформулировать условия, при которых наблюдаемое совпадение может быть интерпретировано как следствие физической гипотезы, а не как численная случайность: (i) диагностировать дискретные выборы постановки, (ii) сформулировать вариационный принцип для отбора $K$ и стабилизации масштаба, (iii) задать соответствие между регуляризованными функциональными детерминантами и наблюдаемой $\alpha(0)$, (iv) выделить независимые предсказания.
\end{abstract}

\section*{Ключевые слова}
Постоянная тонкой структуры; компактификация; линзовые пространства; $\mathbb{RP}^3$; zeta-регуляризация; функциональные детерминанты; стабилизация модулей.

\section{Введение}
\subsection{Мотивация}
Постоянная тонкой структуры $\alpha$ является безразмерной константой, поэтому для неё в принципе возможно ``чисто числовое'' происхождение: значение могло бы быть выражено через безразмерные инварианты некоторой фундаментальной геометрической структуры.
В геометрических и спектральных подходах такой структурой естественно выступает компактный внутренний фактор $K$ в произведении $M=M_4\times K$, поскольку именно компактность переводит спектры эллиптических операторов (Лапласа, Дирака и т.п.) в дискретные и позволяет определить регуляризованные функциональные детерминанты.

\subsection{Проблема статуса результата}
Формулы вида ``$\alpha^{-1}=f(\pi)$'' легко получить постфактум, и поэтому центральной трудностью является не сам факт совпадения чисел, а \emph{предсказательный статус}.
В данной работе мы не утверждаем, что $\alpha^{-1}$ выведено из первых принципов; вместо этого формулируется минимальный набор условий, при которых наблюдаемое численное совпадение может быть физически значимым.
Важно также зафиксировать область утверждений: ниже обсуждается только связка ``$U(1)$-сектор + компактификация на $K=\mathbb{RP}^3\times S^1$ + ζ-регуляризация''; любые более широкие реконструкции параметров Стандартной модели, а также дополнительные числовые/комбинаторные конструкции, не используются как вход и не являются частью доказательной базы данного текста.
Под предсказательностью здесь понимается следующее: (i) выбранная геометрия $K$ и её параметры должны быть следствием вариационного принципа или симметрий, (ii) все дискретные данные постановки (spin-структура, граничные условия, обращение с нулевыми модами, схема вычитания локальных контрчленов) должны быть либо выведены, либо перечислены как независимые физические постулаты с контролем чувствительности результата.

\subsection{Происхождение ``первых принципов'' (границы дедукции)}
В исходной мотивации геометрия $K=\mathbb{RP}^3\times S^1$ возникает как попытка дедукции из структурных требований (внутренний $SU(2)$-сектор, условие реальности/факторизация по центру $\mathbb{Z}_2$, периодичность модулярного потока/KMS).
Важное методологическое различие: (i) вывод самой топологии $K$ из этих предпосылок и (ii) сопоставление наблюдаемой $\alpha(0)$ с конечной частью 1-loop детерминантов на $K$ являются логически разными шагами.
Ниже текст организован так, чтобы (a) явно отделять геометрические/топологические следствия предпосылок от (b) постановочных допущений квантового расчёта и от (c) численных проверок, и тем самым не подменять ``вывод из первых принципов'' постфактум-подгонкой.

\subsection{Идея подхода}
Рабочая гипотеза состоит в том, что низкоэнергетическая (IR) нормировка электромагнитного сектора может быть сопоставлена значению евклидова эффективного действия, вычисленного на внутренней геометрии $K$.
В евклидовой формулировке калибровочно-инвариантная 1-loop структура приводит к комбинациям логарифмов детерминантов (векторный оператор и ghost), которые удобно определять через дзета-регуляризацию:
\begin{equation}
\ln\det \mathcal{O}:= -\zeta_{\mathcal{O}}'(0).
\end{equation}
Таким образом, ключевым объектом становятся спектральные функции $\zeta_{\mathcal{O}}(s)$ и связанные с ними инварианты (включая $\eta$-инвариант и коэффициенты теплового ядра).

\subsection{Центральное утверждение (что именно требуется доказать)}
Ниже центральным является утверждение следующего типа: в фиксированной постановке (паспорт постановки, выбранная схема вычитания локальных контрчленов, правило $\det'$, учёт гармонических мод конечномерным интегралом) калибровочно-инвариантная 1-loop комбинация детерминантов на $K$ задаёт конечную поправку к коэффициенту при $\int_{M_4}F_{\mu\nu}F^{\mu\nu}$ в 4D-эффективном действии, а значит и к $\alpha(0)$ в смысле \eqref{eq:renorm_condition_alpha0}--\eqref{eq:e0_from_Pi}.
Иными словами, требуется доказать (или чётко сформулировать как гипотезу сопоставления), что существует определённая величина $\Delta_K$, выражаемая через одну и ту же калибровочно-инвариантную комбинацию регуляризованных детерминантов на $K=\mathbb{RP}^3\times S^1$, такая что
\begin{equation}
  \frac{1}{e^2(0)}
  =
  \frac{1}{e^2_{\mathrm{4D}}(0)}
  +
  \Delta_K(0;R),
\end{equation}
где $\Delta_K(0;R)$ вычисляется из 1-loop функционального интеграла как конечная (после вычитания локальных контрчленов) часть выражения вида
\begin{equation}
  \Gamma^{(1)}_K
  =
  \frac12\ln\det'\Delta_1
  -\ln\det'\Delta_0
  \;(+\;\text{при необходимости }-\ln\det D\!\!\!\!/\,),
\end{equation}
с пониманием, что гармонические моды (в частности Wilson-line вдоль $S^1$) исключаются из детерминантов и учитываются отдельно.
Весь последующий текст организован вокруг уточнения статуса этого утверждения: какие его части являются стандартной структурой 1-loop QFT (и потому ``жёсткие''), какие сводятся к геометрии/топологии $K$, а какие остаются задачами явного спектрального вычисления.

\subsection{Цель и вклад текста}
Цель работы --- превратить ``численное совпадение'' в проверяемую программу: определить, какие части схемы являются строгими геометрическими фактами, какие --- следствиями общей структуры функционального интеграла, а какие требуют отдельного спектрального вычисления и/или вариационного обоснования.
В частности, формулируются (1) диагностические критерии для выбора $K$, (2) кандидаты на вариационный функционал $\mathcal{F}[K]$, (3) набор сравнительных геометрий для теста, (4) список независимых наблюдательных следствий, необходимых для отделения подгонки от гипотезы.

\subsection{Структура статьи}
Дальнейшее изложение организовано следующим образом.
В разделе 2 перечисляются дискретные выборы постановки и формулируются требования к их выводу.
В разделе 3 описывается вариационный принцип для отбора $K$ и предлагается тестовый набор линзовых пространств.
В разделе 4 обсуждается, как связать $\Gamma_{\mathrm{eff}}$ с $\alpha(0)$ в рамках выбранной схемы регуляризации.
В разделе 5 намечается связь с ренормгруппой и стандартной структурой КЭД.
В разделе 6 приводятся направления для независимых предсказаний.

\subsection{Статус результатов и незакрытые места}
Ниже фиксируем, какие элементы программы \emph{уже} имеют рабочий (частично доказательный) вывод в рамках текста, а какие остаются ключевыми открытыми пунктами.

\begin{itemize}
  \item \textbf{1-loop структура $\Gamma^{(1)}$ (gauge fixing, FP, ghost, разложение Ходжа, $\det'$ и нулевые/гармонические моды).}
  Общая архитектура 1-loop комбинаций детерминантов (векторный оператор $\Delta_1$ и ghost $\Delta_0$, а также фермионный детерминант) в статье уже выписана; однако для журнального уровня требуется вынести в \emph{приложение} компактный, самодостаточный вывод со знаками, нормировками и чётким правилом учёта нулевых и гармонических мод на $K=\mathbb{RP}^3\times S^1$.

  \item \textbf{RG-согласование $\alpha(0)$ vs $\alpha(m_Z)$.}
  В тексте присутствует мост ``ИК-нормировка $\alpha(0)$ $\to$ стандартный ренормгрупповой бег'' на уровне 1-loop формулы; остаётся сформулировать (и явно зафиксировать) схему/пороги, в рамках которых сравнение корректно.

  \item \textbf{Единственность ``геометрического ядра'' $4\pi^3+\pi^2$ против ближайших альтернатив.}
  Аргумент против постфактум-конструирования комбинаций из $\pi$ реализуется как проверка: небольшие вариации целочисленных коэффициентов при $\pi^3,\pi^2$ уводят результат далеко от экспериментального значения.
  Для строгого текста это следует оформить как утверждение внутри фиксированного класса постановок (а не как ``поиск красивой формулы'').

  \item \textbf{Численная реконструкция $R$ из уравнения $\alpha^{-1}(R)=\alpha^{-1}_{\mathrm{exp}}$.}
  В статье уже отмечено, что корень этого уравнения оказывается чрезвычайно близок к $R=1$.
  Однако это \emph{не} заменяет заявленную задачу стабилизации: требуется построение $V_{\mathrm{eff}}(R)$ (или $\Gamma_{\mathrm{eff}}(R)$ на единицу 4D-объёма) и демонстрация устойчивого минимума.

  \item \textbf{Кандидатная финальная формула для $\alpha^{-1}$ и обсуждение $\kappa_{\mathrm{Cas}}=1/24$.}
  В тексте приведён 1D регуляризованный остаток $\kappa_{\mathrm{1D}}=1/24$ и показано, почему в калибровочно-инвариантной комбинации возникает ``половинный'' вклад.
  При этом строгий спектральный вывод того, что \emph{именно эта} константа переносится без схемной неоднозначности в физическую нормировку (после вычитания локальных контрчленов и корректного учёта $\det'$/меры), остаётся отдельной технической задачей.

  \item \textbf{Обоснование выбора $K=\mathbb{RP}^3\times S^1$ в классе альтернатив.}
  ``Быстрый отбор'' кандидатов в классе сферических пространственных форм и линзовых пространств в тексте уже сформулирован (выделение $|\Gamma|=2$ и $p=2$ при требовании дискретного $U(1)$-flat с шагом $\pi$).
  Открытый пункт: замена эвристики вариационным принципом (минимум функционала $\mathcal{F}[K]$) и проверка устойчивости к изменению постановочных данных.

  \item \textbf{Топологический член масштаба $\pi$.}
  В статье подчёркнуто, что это слабейший элемент формулы; необходима либо строгая демонстрация появления универсальной константы из отношения детерминантов в двух плоских секторах ($\mathrm{Hol}=\pm 1$), либо контролируемая замена $\pi$ на другой универсальный инвариант (например, $\log 2$), с пересчётом всей схемы.

  \item \textbf{Контроль схемозависимости через heat-kernel/$a_n$ и коэффициенты при $\log\mu$.}
  На фоне метрики произведения для $\mathbb{RP}^3\times S^1$ и в калибровочно-инвариантной комбинации операторов (Maxwell+ghost, а также возможный фермионный вклад) влияние локальных 4D контрчленов можно существенно ограничить, поскольку $\mu$-зависимость в ζ-регуляризации целиком контролируется величинами вида $\zeta_{\mathcal{O}}(0)$ и, эквивалентно, коэффициентами теплового ядра $a_4(\mathcal{O})$.
  При этом утверждения о занулении конкретных инвариантов конформной аномалии ($W^2$, $E_4$) на $K=\mathbb{RP}^3\times S^1$ требуют отдельной проверки для выбранной метрики и не должны использоваться как неявное допущение.
  Остаётся аккуратно замкнуть этот контроль на конкретной калибровочно-инвариантной комбинации операторов.
\end{itemize}

\subsection{План повышения строгости (с чего начинать)}
Ниже приведён минимальный порядок шагов, который переводит ``совпадение чисел'' в проверяемую постановку с контролем схемозависимости.
Во всех пунктах важно фиксировать 
(i) \\emph{что именно считается наблюдаемым входом} (и на каком масштабе),
(ii) \\emph{что именно выводится} (и в какой схеме),
(iii) \\emph{какие данные постановки являются дискретными} и потому должны быть либо выведены, либо объявлены постулатами.

\paragraph{Целевой формат результата (что должно получиться по итогу).}
Нужно прийти к цепочке утверждений вида
\begin{equation}
\left.\frac{\delta^2\Gamma_{\mathrm{eff}}^{(4)}[A]}{\delta A\,\delta A}\right|_{q^2\to 0}
\Longrightarrow
\frac{1}{e^2(0)}
\Longrightarrow
\alpha^{-1}(0)=\frac{4\pi}{e^2(0)}
\end{equation}
с явным выражением $e^2(0)$ через комбинацию детерминантов на $K$ (и с явным указанием, какие локальные контрчлены вычтены).

\begin{enumerate}
  \item \textbf{Зафиксировать ренормализационное утверждение (IR-нормировка для $\alpha(0)$).}
  Явно записать ренормусловие на 4D-эффективный функционал:
  \begin{equation}\label{eq:renorm_condition_alpha0}
    \Gamma_{\mathrm{eff}}^{(4)}[A]
    \supset
    \frac{1}{4e^2(0)}\int_{M_4} F_{\mu\nu}F^{\mu\nu}
    \qquad (q^2\to 0),
  \end{equation}
  и связать его с 1PI двухточечной функцией фотона. В евклидовом импульсном пространстве вводим
  \begin{equation}\label{eq:gamma2_def}
    \Gamma_{\mathrm{eff}}^{(4)}\supset
    \frac{1}{2}\int\!\frac{d^4q}{(2\pi)^4}\,A_\mu(-q)\,\Pi_{\mu\nu}(q)\,A_\nu(q),
  \end{equation}
  а калибровочная инвариантность задаёт разложение
  \begin{equation}\label{eq:Pi_decomp}
    \Pi_{\mu\nu}(q)=\bigl(q^2\delta_{\mu\nu}-q_\mu q_\nu\bigr)\,\Pi_T(q^2)+\frac{1}{\xi}\,q_\mu q_\nu.
  \end{equation}
  Тогда ИК-нормировка определяется как
  \begin{equation}\label{eq:e0_from_Pi}
    \frac{1}{e^2(0)}:=\lim_{q^2\to 0}\Pi_T(q^2),
    \qquad
    \alpha(0)=\frac{e^2(0)}{4\pi}.
  \end{equation}
  Далее требуется зафиксировать, какая часть $\Pi_T$ относится к ``обычному'' 4D вкладу (с выбранным ренормусловием при масштабе $\mu$), а какая --- к конечной поправке от компактификации на $K$:
  \begin{equation}\label{eq:Pi_split}
    \Pi_T(q^2;\mu)=\Pi_T^{\mathrm{4D}}(q^2;\mu)+\Delta_K(q^2;R,\mu),
  \end{equation}
  где $\Delta_K$ должна быть выражена через одну и ту же калибровочно-инвариантную комбинацию регуляризованных детерминантов (вектор+ghost, с правилом $\det\,'$) на $K$ в фиксированной схеме вычитания локальных контрчленов.
  
  \item \textbf{Перечислить и ``заморозить'' данные постановки (чтобы не было скрытых ручек).}
  Для каждого оператора и поля зафиксировать:
  \begin{itemize}
    \item калибровочную фиксацию (параметр $\xi$) и соответствующие ghost (в абелевом случае --- скалярные грассмановы поля, определяющие $\det\Delta_0$);
    \item точный набор полей в функциональном интеграле (Maxwell+ghost; отдельно оговорить, включается ли материя и на каких расслоениях);
    \item выбор spin-структур на $\mathbb{RP}^3$ и на $S^1$ (в частности, пояснить физический принцип, фиксирующий одну из двух spin-структур на $\mathbb{RP}^3$), поскольку спектр Дирака и детерминанты зависят от этого выбора;
    \item граничные условия/периодичности на $S^1$: для бозонов периодические, для ghost --- те же (чтобы сохранить БРСТ-симметрию);
    \item физическую интерпретацию $S^1$: как пространственной KK-окружности (а не термального KMS-круга), что фиксирует выбор spin-структуры/периодичности для фермионов и исключает трактовку этого выбора как свободной ``ручки'';
    \item топологические данные на $K$: класс $U(1)$-расслоения и допустимые дискретные голономии (например, тривиальный класс Черна при возможном нетривиальном плоском секторе из $\pi_1(\mathbb{RP}^3)=\mathbb{Z}_2$);
    \item правило обращения с нулевыми модами и нормировки меры: $\det\,'$ (исключение нулевых собственных значений операторов) + единое соглашение о делении на объём калибровочной группы;
    \item отдельное правило для гармонических мод: поскольку $b_1(\mathbb{RP}^3\times S^1)=1$, у 1-форм имеется гармоническая мода вдоль $S^1$; требуется явно указать, интегрируем ли мы по ней как по физическому модулю, или исключаем её как калибровочную/сводим к сумме по дискретным голономиям;
    \item схему вычитания локальных контрчленов (какие коэффициенты теплового ядра вычитаются и почему), причём эта схема считается частью определения величины, сопоставляемой с $\alpha(0)$.
  \end{itemize}
  Результат этого пункта должен быть одним фиксированным ``паспортом постановки'', к которому можно ссылаться во всех дальнейших вычислениях.

  \item \textbf{Зафиксировать нормировку геометрии и смысл параметра $R$.}
  Дать однозначное определение эталонной метрики $g_0$ на $K$ и масштабирования $g(R)=R^2 g_0$.
  Отдельно выписать, какие геометрические величины дают коэффициенты $4\pi^3,\pi^2,\pi$ в $S_{\mathrm{geo}}$ (объёмы, периоды, фундаментальные области голономий), и при каких соглашениях (кривизна ``единичной'' сферы, длина окружности, нормировка калибровочного поля) эти коэффициенты меняются.

  \item \textbf{Заменить ``геометрическое ядро'' спектральным вычислением.}
  Для каждого вклада в $S_{\mathrm{geo}}$ выписать соответствующий оператор (после gauge fixing и разложения Ходжа) и показать, что требуемая комбинация возникает из
  \begin{equation}
    \ln\det \mathcal{O}:= -\zeta_{\mathcal{O}}'(0)
    \qquad\text{(и/или из коэффициентов теплового ядра),}
  \end{equation}
  именно на $K=\mathbb{RP}^3\times S^1$ и при зафиксированных данных постановки из пункта 2.

  \item \textbf{Стабилизация $R$ как динамическая задача, а не как ``корень уравнения''.}
  Построить $\Gamma_{\mathrm{eff}}(R)$ (на единицу 4D-объёма) из тех же детерминантов и сформулировать условия устойчивого минимума
  $$\left.\partial_R \Gamma_{\mathrm{eff}}\right|_{R=1}=0,\qquad \left.\partial_R^2 \Gamma_{\mathrm{eff}}\right|_{R=1}>0,$$
  контролируя зависимость от схемы вычитания из пункта 2.

  \item \textbf{Тест на уникальность геометрии (внутренний ``кросс-чек'' без подгонки).}
  На ограниченном, но явном классе конкурентов (например, $L(p,q)\times S^1$ при нескольких $p$) повторить вычисление ключевой комбинации детерминантов и проверить, что выбор $\mathbb{RP}^3\times S^1$ выделяется при \\emph{тех же} дискретных правилах из пункта 2.

  \item \textbf{Независимые предсказания/фальсификация.}
  Сформулировать хотя бы одно следствие, зависящее от тех же фиксированных данных постановки, но \\emph{не} использующее $\alpha(0)$ как вход (например, наличие/масштаб масс топологических дефектов $\sim 1/R$, или фиксированная структура дискретных голономий и связанных селекционных правил).
\end{enumerate}

\section{Диагностика дискретных выборов}
Цель этого раздела --- перечислить все места, где в постановке присутствуют дискретные (и потому потенциально ``ручные'') выборы, и явно сформулировать, что именно нужно вывести, чтобы переход от совпадения чисел к предсказательной гипотезе был обоснован.

\subsection{Паспорт постановки (фиксируемые данные)}
Чтобы дальнейшие вычисления имели предсказательный статус, удобно один раз зафиксировать ``паспорт постановки'' --- набор соглашений и дискретных данных, на которые дальше можно ссылаться.

\paragraph{Геометрия и масштаб.}
Внутреннее пространство фиксируется как $K=\mathbb{RP}^3\times S^1$ с эталонной метрикой $g_0$ и масштабированием $g(R)=R^2 g_0$.

\paragraph{Калибровочный сектор.}
Рассматривается абелев $U(1)$-потенциал $A$ с действием Максвелла; калибровочная фиксация берётся в $\xi$-калибровке.
Функциональный интеграл включает соответствующие ghost поля.

\paragraph{Операторы и детерминанты (1-loop).}
После gauge fixing и разложения Ходжа вклад 1-loop должен быть записан в виде калибровочно-инвариантной комбинации детерминантов
(векторный оператор на 1-формах и скалярный ghost-оператор на 0-формах) с дзета-регуляризацией
$\ln\det\mathcal{O}:=-\zeta_{\mathcal{O}}'(0)$.

\paragraph{Нулевые моды и объём калибровочной группы.}
Нулевые моды (в том числе глобальные калибровочные преобразования/гармонические формы при наличии) исключаются правилом $\det\,'$;
нормировка по объёму калибровочной группы фиксируется одним и тем же соглашением во всех сравнениях геометрий.

\paragraph{Топологические данные и голономии.}
Фиксируется класс $U(1)$-связи на $K$ (включая допустимые дискретные голономии, если они используются в аргументе).
Если в дальнейшем привлекаются спиноры, отдельно фиксируется spin-структура и соответствующее расслоение; в частности, на $\mathbb{RP}^3$ необходимо зафиксировать одну из двух spin-структур.
В данной работе принимается (как часть постановки) выбор продолжаемой (bounding) spin-структуры на $\mathbb{RP}^3$ и интерпретация окружности $S^1$ как пространственной KK-окружности, что приводит к периодичности фермионов по $S^1$.

\paragraph{Ренормусловие и смысл $\alpha(0)$ (что считается наблюдаемым).}
Величина, сопоставляемая с $\alpha(0)$, определяется через коэффициент при $\int_{M_4}F_{\mu\nu}F^{\mu\nu}$ в 4D-эффективном действии (см. \eqref{eq:renorm_condition_alpha0}--\eqref{eq:e0_from_Pi}).
Соответственно, геометрическая/спектральная часть постановки должна поставлять \emph{конечную} поправку к этому коэффициенту, причём именно в одной фиксированной калибровочно-инвариантной комбинации детерминантов (Maxwell+ghost, с правилом $\det'$ и с выделением гармонических мод в конечномерный интеграл).

\paragraph{Схема вычитания локальных контрчленов.}
Фиксируется, какие локальные вклады (коэффициенты теплового ядра) вычитаются как 4D-контрчлены и какие остаются как конечные, геометрически нетривиальные остатки.
Именно эта схема далее считается частью определения величины, сопоставляемой с $\alpha(0)$.

\subsection{Сводная таблица: выборы, проблемы и требуемые доказательства}
\begin{center}
\renewcommand{\arraystretch}{1.25}
\begin{tabular}{|p{0.22\linewidth}|p{0.23\linewidth}|p{0.25\linewidth}|p{0.24\linewidth}|}
\hline
\textbf{Выбор} & \textbf{Текущее обоснование} & \textbf{Проблема} & \textbf{Что надо доказать} \\
\hline
$R=1$ & ``планковские единицы'' & калибровка единиц $\neq$ минимум модуля & стабилизация: минимум $V_{\mathrm{eff}}(R)$ \\
\hline
$K=\mathbb{RP}^3\times S^1$ & ``минимальность'' & не вариационный критерий & минимум $\mathcal{F}[K]$ на классе $S^3/\Gamma$ \\
\hline
$4\pi^3,\pi^2$ & ``два сектора'' & не указан оператор/вес & выписать операторы/BC и вычислить $\zeta'(0)$ \\
\hline
$\kappa_{\mathrm{Cas}}=1/24$ & ``универсальность'' & без явного спектрального вывода & вывод из спектра/вычитания (вектор+ghost, $\det'$) \\
\hline
$\delta_{BB}\sim 1/(\pi^4 S_{\mathrm{geo}}^2)$ & ``2-й порядок'' & не указан механизм & кумулянта/2-loop/смешивание, контроль размерностей \\
\hline
\end{tabular}
\end{center}

\paragraph{Текущее состояние (внутренний план развёртывания вывода).}
Ниже мы фиксируем, какие пункты уже имеют рабочий (хотя и не всегда окончательно строгий) вывод и что именно требуется переписать в виде самодостаточных лемм/теорем.
\begin{itemize}
  \item $R=1$: численно проверена сильная чувствительность формулы $\alpha^{-1}(R)$ к масштабу и установлен корень $R\approx 1$; остаётся довести до вариационного утверждения о минимуме $V_{\mathrm{eff}}(R)$.
  \item Выбор $K=\mathbb{RP}^3\times S^1$: показано, что при требованиях (i) нетривиальная $\pi_1$ без непрерывных модулей плоских $U(1)$-связностей и (ii) «шаг» голономии порядка $\pi$ в линзовом классе, выделяется случай $\pi_1=\mathbb{Z}_2$, то есть $\mathbb{RP}^3$.
  \item Вклады $4\pi^3$ и $\pi^2$: разобрана нормировка нулевых мод в KK-редукции и роль накрытия $S^3\to\mathbb{RP}^3$ для спинорного сектора.
  \item Коэффициент $1/24$: выведен как одномерный регуляризованный остаток, согласованный со структурой 1-loop комбинации Maxwell+ghost после gauge fixing и разложения Ходжа.
  \item РГ-интерпретация: аргументировано, что совпадение относится к $\alpha(0)$ (IR), а переход к $\alpha(m_Z)$ задаётся стандартным running.
\end{itemize}


\subsection{Масштаб \texorpdfstring{$R$}{R} и стабилизация модуля}
Фиксация $\hbar=c=G=1$ задаёт единицы измерения, но не решает задачу стабилизации радиуса компактификации.
Поэтому требование ``$R=1$'' должно быть заменено утверждением вида: \emph{эффективный потенциал имеет устойчивый минимум при $R=1$}.

\paragraph{Текущее состояние (что уже можно записать в статье).}
Если восстановить масштабную зависимость геометрических членов при общем радиусе $R$ (в планковских единицах), то
\begin{equation}
S_{\mathrm{geo}}(R)=4\pi^3 R^4+\pi^2 R^3,
\qquad
\alpha^{-1}(0)=\frac{\mathrm{Vol}(K(R))}{g_5^2}+\Delta_{\mathrm{1\text{-}loop}}^{\mathrm{finite}}(R)+\delta_{\mathrm{top}}(R).
\end{equation}

\paragraph{Почему здесь \texorpdfstring{$R$}{R} и \texorpdfstring{$\pi$}{pi} связаны жёстко (а не являются независимыми ``ручками'').}
Ни $R$, ни появление $\pi$ в $S_{\mathrm{geo}}$ не вводятся как независимые числовые параметры: выбирается \emph{эталонная} метрика $g_0$ на $\mathbb{RP}^3\times S^1$ (и на накрытии $S^3\times S^1$ там, где это необходимо) с фиксированной геометрической нормировкой при $R=1$ (например, секционная кривизна сферического фактора равна 1), а затем рассматривается однопараметрическое масштабирование
\begin{equation}
 g(R)=R^2 g_0.
\end{equation}
Тогда все объёмные/локальные геометрические вклады неизбежно приобретают степени $R$ (в частности, $\mathrm{Vol}(\mathbb{RP}^3)\propto R^3$, $\mathrm{Vol}(S^1)\propto R$), тогда как числовые коэффициенты вида $2\pi$, $\pi^2$, $4\pi^3$ фиксируются уже на уровне $g_0$ стандартными формулами для объёмов единичных сферических форм и периодом больших калибровочных преобразований $U(1)$, задающим фундаментальную область для Wilson-линии длины $2\pi$.
В этом смысле корректная постановка задачи состоит не в ``настройке'' $\pi$, а в проверке того, что выбранная нормировка $g_0$ и вариационный отбор $R$ (минимум $V_{\mathrm{eff}}(R)$) совместимы.

Численные проверки, основанные на подстановке в конкретную «закрытую» формулу $\alpha^{-1}(R)=\alpha^{-1}_{\mathrm{exp}}$, не могут подменять стабилизацию: без явного вычисления $\Delta_{\mathrm{1\text{-}loop}}^{\mathrm{finite}}(R)$ и формы $\delta_{\mathrm{top}}(R)$ уравнение для $R$ не является физически осмысленным.
Требуемый динамический результат остаётся тем же: нужно построить $V_{\mathrm{eff}}(R)$ (или $\Gamma_{\mathrm{eff}}(R)$ на единицу 4D-объёма) и показать существование устойчивого минимума.


\paragraph{Что именно считать $R$?}
В данной постановке удобно вводить $R$ как общий масштаб метрики на сферическом факторе (и, при необходимости, на окружности):
\begin{equation}
K(R)=\bigl(\mathbb{RP}^3,\ g_{\mathbb{RP}^3}(R)\bigr)\times \bigl(S^1,\ g_{S^1}(R)\bigr),
\end{equation}
где $\mathrm{Vol}(\mathbb{RP}^3)\propto R^3$, $\mathrm{Vol}(S^1)\propto R$.

\paragraph{Кандидат на $V_{\mathrm{eff}}(R)$.}
В простейшей версии $V_{\mathrm{eff}}(R)$ строится из регуляризованных вакуумных энергий (или из евклидова эффективного действия на единицу 4D-объёма):
\begin{equation}
V_{\mathrm{eff}}(R)=V_{\mathrm{cl}}(R)+V_{\mathrm{1\text{-}loop}}(R)+\cdots,
\end{equation}
где $V_{\mathrm{1\text{-}loop}}$ выражается через детерминанты соответствующих операторов на $K(R)$.
В терминах дзета-регуляризации типичный вклад имеет вид
\begin{equation}
\Gamma_{\mathrm{1\text{-}loop}}(R)=\pm\tfrac12\,\ln\det \mathcal{O}(R)=\mp\tfrac12\,\zeta_{\mathcal{O}(R)}'(0),
\end{equation}
с учётом $\det'$ (исключение нулевых мод) и калибровочно-инвариантных комбинаций.

\paragraph{Критерий стабилизации.}
Необходимо показать
\begin{equation}
\left.\partial_R V_{\mathrm{eff}}\right|_{R=1}=0,
\qquad
\left.\partial_R^2 V_{\mathrm{eff}}\right|_{R=1}>0,
\end{equation}
и, кроме того, что минимум не зависит от произвольного выбора схемы вычитания локальных контрчленов (или что схема фиксируется независимым принципом).

\subsection{Выбор \texorpdfstring{$K=\mathbb{RP}^3\times S^1$}{K=RP3 x S1} вместо других компактификаций}
Тезис ``минимальность'' должен быть переформулирован как задача минимума функционала.
В классе сферических 3-форм $S^3/\Gamma$ естественно сравнивать линзовые пространства $L(p,q)$.

\paragraph{Требуемое утверждение.}
Нужно доказать (или опровергнуть) существование такого функционала $\mathcal{F}[K]$, что
\begin{equation}
\mathcal{F}[\mathbb{RP}^3\times S^1]\le \mathcal{F}[L(p,q)\times S^1]
\end{equation}
для заданного набора $(p,q)$ (и, желательно, для всех допустимых $p$ в выбранном классе), при фиксированных физических ограничениях: spin-структура, нетривиальная $\pi_1$ и постоянная положительная кривизна.


\paragraph{Быстрый отбор в классе сферических форм (что уже следует без тонких вычислений).}
Рассмотрим естественный класс 3-мерных сферических пространственных форм $M^3=S^3/\Gamma$ (свободное действие конечной группы $\Gamma\subset SO(4)$) и внутренние пространства вида $K=M^3\times S^1$.

\begin{enumerate}
  \item \textbf{Фиксация $Z_A=\pi^2$ выделяет $\Gamma=\mathbb{Z}_2$.}
  Поскольку $\mathrm{Vol}(S^3)=2\pi^2$, имеем
  \begin{equation}
  \mathrm{Vol}(M^3)=\frac{\mathrm{Vol}(S^3)}{|\Gamma|}=\frac{2\pi^2}{|\Gamma|}.
  \end{equation}
  Требование $\mathrm{Vol}(M^3)=\pi^2$ эквивалентно $|\Gamma|=2$, то есть в этом классе единственный вариант: $M^3=\mathbb{RP}^3=L(2,1)$.

  \item \textbf{Дискретность $U(1)$-flat и масштаб $\pi$ также выделяют $\pi_1=\mathbb{Z}_2$.}
  Для линзового пространства $L(p,1)$ имеем $\pi_1\cong\mathbb{Z}_p$, а пространство модулей плоских $U(1)$-связностей с точностью до калибровки
  \begin{equation}
  \mathcal{M}_{\mathrm{flat}}\bigl(L(p,1),U(1)\bigr)\cong \mathrm{Hom}(\mathbb{Z}_p,U(1))\cong \{e^{2\pi i k/p}\}_{k=0}^{p-1}.
  \end{equation}
  Это $p$ дискретных вакуумов с характерным «шагом» $2\pi/p$; чтобы получить ровно два вакуума и масштаб $\pi$, требуется $p=2$.

  \item \textbf{Исключение непрерывных модулей.}
  Если $b_1(M^3)>0$ (например, для $T^3$ или $S^2\times S^1$), то $\mathcal{M}_{\mathrm{flat}}(M^3,U(1))$ содержит непрерывные компоненты, и в функциональном интеграле возникает дополнительная (схемозависимая) нормировка меры по модулю.
  Для $\mathbb{RP}^3$ имеем $b_1(\mathbb{RP}^3)=0$.
\end{enumerate}

Этот «быстрый отбор» не заменяет вариационный принцип, но показывает, что в наиболее естественном классе геометрий $M^3\times S^1$ структура $S_{\mathrm{geo}}=4\pi^3+\pi^2+\pi$ без введения дополнительных безразмерных параметров фактически вынуждает выбор $M^3=\mathbb{RP}^3$.
Практическая проверка (в том числе перебор близких альтернативных комбинаций коэффициентов при $\pi^3,\pi^2,\pi$ и оценка расхождения с $\alpha^{-1}_{\mathrm{exp}}$) может быть выполнена численно; в тексте ниже она используется лишь как вспомогательный тест против постфактум-конструирования.

\paragraph{Почему одной систолы недостаточно.}
Систола является полезным геометрическим инвариантом, но она не учитывает спектральную информацию и не контролирует калибровочно-инвариантную комбинацию детерминантов.
Поэтому даже если $\mathbb{RP}^3$ ``выгодна'' по $\mathrm{sys}$, это ещё не означает минимума полной вакуумной энергии.

\subsection{Идентификация вкладов $4\pi^3$, $\pi^2$, $\pi$}
Ключевая слабость текущей ``суммы трёх чисел'' --- отсутствие явной карты ``оператор $\to$ детерминант $\to$ геометрический инвариант''.
Для журнальной версии нужно ввести конкретный набор полей и операторов на $K$.

\paragraph{Текущее состояние (развёрнутый вывод KK-нормировок).}
Здесь требуется два независимых утверждения: (i) как в 5D Maxwell-секторе возникает $\pi^2$, и (ii) почему в фермионном секторе появляется именно $4\pi^3$.

\paragraph{(i) Калибровочный сектор и $Z_A=\pi^2$.}
Начинаем с 5D действия Maxwell на $M_4\times (\mathbb{RP}^3\times S^1)$:
\begin{equation}
S_A=\frac{1}{4g_5^2}\int_{M_4} d^4x\int_{\mathbb{RP}^3\times S^1} d\mathrm{vol}\; F_{MN}F^{MN}.
\end{equation}
Для нулевой моды по внутренним координатам получаем 4D действие
\begin{equation}
S_A=\frac{\mathrm{Vol}(\mathbb{RP}^3\times S^1)}{4g_5^2}\int_{M_4} d^4x\;F_{\mu\nu}F^{\mu\nu}
\equiv \frac{1}{4g_4^2}\int_{M_4} d^4x\;F_{\mu\nu}F^{\mu\nu},
\end{equation}
то есть
\begin{equation}
\frac{1}{g_4^2}=\frac{\mathrm{Vol}(\mathbb{RP}^3\times S^1)}{g_5^2}.
\end{equation}
Далее используется стандартная нормировка $U(1)$-связности: голономия вдоль окружности задаётся переменной Wilson-line
$\theta:=\int_{S^1}A\in[0,2\pi)$.
Большие калибровочные преобразования сдвигают $\theta\mapsto\theta+2\pi n$, поэтому физический модуль плоской компоненты $A_{S^1}$ интегрируется по фундаментальной области длины $2\pi$ и тем самым фиксирует, что единственная безразмерная комбинация в 5D Maxwell-секторе имеет вид $g_5^2/\mathrm{Vol}(S^1)$.
Если не вводить дополнительного безразмерного параметра (чтобы не потерять предсказательность), принимается нормировка
\begin{equation}
\frac{g_5^2}{\mathrm{Vol}(S^1)}=1
\qquad\Longleftrightarrow\qquad
g_5^2=\mathrm{Vol}(S^1).
\end{equation}
Тогда при радиусе $R=1$ имеем $\mathrm{Vol}(S^1)=2\pi$ и $\mathrm{Vol}(\mathbb{RP}^3\times S^1)=\pi^2\cdot 2\pi$, откуда
\begin{equation}
\frac{1}{g_4^2}=\frac{\pi^2\cdot 2\pi}{2\pi}=\pi^2.
\end{equation}

\paragraph{(ii) Фермионный сектор и $Z_\Psi=4\pi^3$.}
Ключевое отличие спинорного сектора состоит в том, что при фиксированной spin-структуре на $\mathbb{RP}^3=S^3/\mathbb{Z}_2$ спектральные моды спинора соответствуют нечётным гармоникам относительно антиподного отображения.
Такие моды естественно нормировать на накрытии $S^3$, а не на факторе, поэтому нормировочный интеграл для спинора берётся по $S^3\times S^1$:
\begin{equation}
Z_\Psi=\mathrm{Vol}(S^3\times S^1)=\mathrm{Vol}(S^3)\,\mathrm{Vol}(S^1)=2\pi^2\cdot 2\pi=4\pi^3.
\end{equation}
Именно эта разница в «видимой» геометрии (фактор vs накрытие) и порождает раздельные вклады $4\pi^3$ и $\pi^2$ в суммарном $S_{\mathrm{geo}}$.



\paragraph{Минимальный набор операторов (пример).}
Для $U(1)$-калибровочного сектора стандартна комбинация
\begin{equation}
\Gamma^{(1)}_{\mathrm{gauge}}=\tfrac12\ln\det'\Delta_1-\ln\det'\Delta_0,
\end{equation}
где $\Delta_1$ --- лапласиан на 1-формах, $\Delta_0$ --- лапласиан на скалярах (ghost).
Для фермионов (при заданной spin-структуре и BC по $S^1$) возникает
\begin{equation}
\Gamma^{(1)}_{\mathrm{ferm}}=-\ln\det D,
\end{equation}
с выбором спектральной ветви и правилом обращения с нулевыми модами.

\paragraph{Воспроизводимость (что именно должно быть самодостаточно в статье).}
Для журнального статуса необходимо, чтобы фиксация знаков, нормировок и обращения с нулевыми/гармоническими модами (включая обсуждение $b_0=1$ и $b_1=1$ для $K=\mathbb{RP}^3\times S^1$) была приведена прямо в тексте (или в приложении) без опоры на внешние материалы.
Аналогично, sanity-проверка ренормгруппового бега должна быть оформлена как короткая аналитическая проверка (с порогами), а не как ссылка на внешний код.

\paragraph{Почему объём может появляться.}
В разложении теплового ядра
\begin{equation}
\mathrm{Tr}\,e^{-t\mathcal{O}}\sim \sum_{n\ge 0} a_n(\mathcal{O})\,t^{(n-d)/2},
\end{equation}
ведущий коэффициент $a_0$ пропорционален $\mathrm{Vol}(K)$.
Однако сам по себе $a_0$ отвечает локальной (вейлевской) части и обычно подлежит вычитанию/перенормировке.
Поэтому требуется чётко указать, \emph{какой} именно (локальный или нелокальный) остаток фиксируется как физический и почему он в выбранной схеме даёт именно $4\pi^3$ и $\pi^2$ с весами 1.

\paragraph{Топологический вклад $\pi$ (самое слабое место и критерий проверки).}
В формуле $S_{\mathrm{geo}}=4\pi^3+\pi^2+\pi$ именно слагаемое $\pi$ является наименее строго обоснованным: оно не следует напрямую ни из объёма (как $\pi^2$), ни из ``накрытия'' (как $4\pi^3$).
Поэтому мы фиксируем его не как ``доказанный факт'', а как \emph{проверяемую гипотезу о топологическом вкладе}.

\paragraph{Гипотеза $\mathbf{H_\pi}$.}
Существует калибровочно-инвариантная величина в 1-loop эффективном действии $\Gamma_{\mathrm{eff}}$ на $K=\mathbb{RP}^3\times S^1$, которая в выбранной схеме редукции и нормировки даёт аддитивный вклад масштаба порядка
\begin{equation}
\Delta_{\pi}\Gamma\sim L_{\mathrm{sys}}(\mathbb{RP}^3)=\pi\quad (R=1).
\end{equation}
Смысл гипотезы: дискретная топология $\pi_1(\mathbb{RP}^3)=\mathbb{Z}_2$ (две плоские $U(1)$-связности) должна проявляться не только как ``наличие двух секторов'', но и давать конкретный вклад в конечную часть нормировки.

\paragraph{Почему это реально слабое место (и где потенциальная ошибка).}
Наиболее распространённые спектральные топологические инварианты в данной постановке не дают $\pi$:
для стандартной метрики (и выбранной spin-структуры) $\eta$-инвариант оказывается равным нулю, а логарифм аналитического (Ray--Singer) торсиона для $L(2,1)$ даёт $\log T_{RS}=-2\log 2$.
Это означает: \emph{если искомый вклад действительно топологический и спектральный}, то он не должен быть отождествлён напрямую с $\eta$ или торсионом.

\paragraph{Критерий фальсификации (что именно считаем ``закрытием'' вопроса).}
Вопрос о $\pi$ считаем закрытым только если выполнено одно из двух:
\begin{enumerate}
  \item \textbf{Строгий вывод:} предъявлен конкретный сектор функционального интеграла (какое поле/какие BC/какая плоская связь) и вычислена конечная часть соответствующей калибровочно-инвариантной комбинации ζ-детерминантов, дающая вклад, пропорциональный $\pi$.

  \item \textbf{Контролируемая замена:} показано, что корректный спектральный инвариант даёт другой универсальный вклад (например, $\log 2$ или $2\log 2$), и тогда слагаемое $\pi$ в $S_{\mathrm{geo}}$ должно быть заменено на этот инвариант, после чего вся модель переоценивается (включая численную подгонку и выводы о $R\approx1$).
\end{enumerate}

\paragraph{Минимальный план вычисления для \texorpdfstring{$\pi_1=\mathbb{Z}_2$}{pi1=Z2} (без расширения модели).}
Чтобы не оставлять $\pi$ ``по вкусу'', фиксируем конкретную вычислительную цель.
Для $U(1)$ на $\mathbb{RP}^3$ имеется два дискретных класса плоских связностей (тривиальная и нетривиальная, $\mathrm{Hol}=\pm1$).
Нужно:
\begin{enumerate}
  \item выбрать оператор, реально входящий в калибровочно-инвариантную 1-loop комбинацию (например, $\Delta_1/\Delta_0$ или оператор Дирака с $U(1)$-связностью);
  \item вычислить отношение детерминантов в двух плоских секторах,
  \begin{equation}
  \frac{\det'\mathcal{O}[\mathrm{Hol}=-1]}{\det'\mathcal{O}[\mathrm{Hol}=+1]},
  \end{equation}
  и извлечь из него вклад в конечную часть нормировки (после однозначно зафиксированного вычитания локальных контрчленов);
  \item проверить, какой универсальный ``топологический'' масштаб возникает: $\pi$, $\log 2$, либо иная константа.
\end{enumerate}
Именно этот расчёт (а не геометрическая эвристика) должен решить судьбу $\pi$.

\paragraph{Воспроизводимый 1D-тест (круг с твистом).}
Как минимальную sanity-проверку того, что отношение детерминантов действительно может порождать логарифмы фундаментальных геометрических констант, рассмотрим скалярный лапласиан на окружности $S^1$ радиуса $R$ с квазипериодическими граничными условиями
\begin{equation}
  f(x+2\pi R)=e^{2\pi i\phi}f(x),\qquad \phi\in[0,1).
\end{equation}
Тогда спектр равен $\lambda_n=((n+\phi)/R)^2$, $n\in\mathbb{Z}$.
Для тривиального сектора $\phi=0$ имеется нулевая мода, поэтому используем $\det'$ (исключая $n=0$).
В твистованном секторе $\phi=\tfrac12$ (антипериодичность, что в абелевом контексте соответствует простейшей модели $\mathrm{Hol}=-1$) получаем ζ-регуляризованные детерминанты
\begin{equation}
  \ln\det\Delta_{S^1,\,\phi=1/2}=2\ln\bigl(2\sin(\pi/2)\bigr)=2\ln 2,
  \qquad
  \ln\det'\Delta_{S^1,\,\phi=0}=2\ln(2\pi R).
\end{equation}
Следовательно,
\begin{equation}\label{eq:S1_det_ratio_pi}
  \ln\frac{\det\Delta_{S^1,\,\phi=1/2}}{\det'\Delta_{S^1,\,\phi=0}}
  =
  2\ln\frac{2}{2\pi R}
  =
  -2\ln(\pi R).
\end{equation}
В частности, при $R=1$ имеем
\begin{equation}
  \ln\frac{\det\Delta_{S^1,\,\phi=1/2}}{\det'\Delta_{S^1,\,\phi=0}}=-2\ln\pi.
\end{equation}
Численная проверка формул и конвенций (включая правило $\det'$) реализована в скрипте \texttt{scripts/zeta\_det\_circle\_holonomy.py}.

\paragraph{Toy Maxwell+ghost: разность «секторов» на \texorpdfstring{$S^3$}{S3}.}
Чтобы отделить эффект дискретного сектора от объёмной (локальной) части, полезно сравнить ``нескрученный'' тепловой след на $S^3$ с ``скрученным'' следом $\mathrm{Tr}(g\,e^{-t\Delta})$, моделирующим вклад $\mathbb{Z}_2$-сектора (антипод $g$).
В скрипте \texttt{scripts/pi\_sector\_maxwell\_ghost\_holonomy\_scan\_S3.py} реализована toy-комбинация Maxwell+ghost
\begin{equation}
  \Delta\Gamma_{\mathrm{toy}}:=\tfrac12 I_1-I_0,
  \qquad
  I_p:=-\int d(\ln t)\,K^{(p)}(t),
\end{equation}
в двух вариантах:
(+) обычный след на $S^3$ и (-) скрученный след с вставкой $g$.
При параметрах $R=1$, $t\in[10^{-3},50]$, $n_{\max}=8000$ получается
\begin{equation}
  \Delta\Gamma_{+}\approx 8.24042308004\times 10^{1},
  \qquad
  \Delta\Gamma_{-}\approx 1.11395063327,
\end{equation}
и потому формальная разность
\begin{equation}
  \Delta\Gamma_{-}-\Delta\Gamma_{+}\approx -8.12902801671\times 10^{1}
\end{equation}
устойчива к увеличению $n_{\max}$.
Смысл этого теста следующий: ``скрученный'' сектор естественно даёт конечную константу порядка 1, тогда как ``нескрученный'' объект содержит большую локальную (объёмную) часть.
Следовательно, физически осмысленная ``разность секторов'' может быть определена только после явного вычитания локальных контрчленов (то есть после выделения конечной, нелокальной части), и именно эта конечная часть должна сравниваться с кандидатами вроде $\pi$ или $\log 2$.

\paragraph{Toy Maxwell+ghost на \texorpdfstring{$K=\mathbb{RP}^3\times S^1$}{K=RP3xS1}: два Hol-сектора.}
Следующий шаг --- подняться с 3D-теста на модельное внутреннее пространство $K=\mathbb{RP}^3\times S^1$.
В скрипте \texttt{scripts/pi\_sector\_maxwell\_ghost\_RP3xS1\_holonomy.py} два плоских сектора $\mathrm{Hol}=\pm1$ на $\mathbb{RP}^3$ реализуются через проектор на накрытии $S^3$:
\begin{equation}
  \mathrm{Tr}_{\mathbb{RP}^3,\,\rho}(e^{-t\Delta})
  =\frac12\Bigl(\mathrm{Tr}_{S^3}(e^{-t\Delta})+\rho(g)\,\mathrm{Tr}_{S^3}(g\,e^{-t\Delta})\Bigr),
  \qquad \rho(g)=\pm1.
\end{equation}
Фактор $S^1$ учитывается KK-суммой $\Theta(t)=\sum_{m\in\mathbb{Z}}e^{-t(m/R_1)^2}$, а гармоническая 1-форма вдоль $S^1$ исключается правилом $\det'$ (в toy-реализации как вычитание единицы из теплового следа 1-форм).
Локальная (4D) часть вычитается в toy-схеме через подгонку $K_{\mathrm{loc}}(t)=A t^{-2}+B t^{-1}$ в малом окне по $t$.
При параметрах $R_3=R_1=1$, $t\in[10^{-3},50]$, fit-окно $t\in[2\times10^{-3},2\times10^{-2}]$ и срезах $n_{\max}=8000$, $m_{\max}=4000$ получено
\begin{equation}
  \Delta\Gamma^{\mathrm{ren}}(\mathrm{Hol}=+1)\approx -12.1236048194,
  \qquad
  \Delta\Gamma^{\mathrm{ren}}(\mathrm{Hol}=-1)\approx 4.14856084256,
\end{equation}
и, следовательно,
\begin{equation}
  \Delta\Gamma^{\mathrm{ren}}_{\mathrm{top}}
  :=
  \Delta\Gamma^{\mathrm{ren}}(\mathrm{Hol}=-1)-\Delta\Gamma^{\mathrm{ren}}(\mathrm{Hol}=+1)
  \approx 16.272165662.
\end{equation}
Этот toy-результат демонстрирует, что даже на $\mathbb{RP}^3\times S^1$ (при грубой 4D-перенормировке) разность Hol-секторов не выделяет автоматически $\pi$ и чувствительна к выбранной схеме вычитания; тем самым подтверждается необходимость полного 4D background-field вывода коэффициента при $\int F^2$.

\subsection{Происхождение коэффициента $1/24$}
Число $1/24$ напоминает стандартные значения дзета-функции ($\zeta_R(-1)=-1/12$) и возникает в одномерных касимировских остатках.
Для журнального статуса необходимо воспроизвести этот коэффициент из конкретной спектральной суммы.

\paragraph{Текущее состояние (кандидатный строгий вывод).}
Здесь требуется вывести безразмерную константу $\kappa_{\mathrm{Cas}}$ как конечную часть одномерной спектральной суммы, совместимую с калибровочно-инвариантной 1-loop комбинацией детерминантов.

\paragraph{Шаг 1: чисто одномерный остаток.}
Рассмотрим Abel-суммирование (эквивалентное ζ-регуляризации для полиномиальных расходимостей)
\begin{equation}
S(t):=\sum_{n\ge 1} n\,e^{-tn},\qquad t>0.
\end{equation}
Эта сумма вычисляется в элементарных функциях:
$S(t)=\frac{e^{-t}}{(1-e^{-t})^2}$, и при $t\to 0^+$ имеет асимптотику Бернулли
\begin{equation}
S(t)=\frac{1}{t^2}-\frac{1}{12}+O(t^2).
\end{equation}
Поэтому существует конечный регуляризованный остаток
\begin{equation}
\kappa_{\mathrm{1D}}:=-\frac12\lim_{t\to 0^+}\Bigl(S(t)-\frac{1}{t^2}\Bigr)=\frac{1}{24},
\end{equation}
что эквивалентно записи $\kappa_{\mathrm{1D}}=-\zeta_R(-1)/2$ при $\zeta_R(-1)=-1/12$.

\paragraph{Шаг 2: почему появляется множитель $1/2$ в QED.}
Для $U(1)$-сектора после gauge fixing и введения ghost-полей 1-loop часть эффективного действия имеет структуру
\begin{equation}
\Gamma^{(1)}_{\mathrm{gauge}}=\frac12\log\det'(\Delta_1)-\log\det'(\Delta_0).
\end{equation}
Разложение Ходжа для 1-форм $A=A_{\perp}+d\phi$ даёт факторизацию спектра на продольные и поперечные моды, так что продольная часть $\det'(\Delta_0)$ сокращается с ghost-произведением и остаётся ``половинная'' комбинация
\begin{equation}
\Gamma^{(1)}_{\mathrm{gauge}}=\frac12\log\det'(\Delta_{1,\perp})-\frac12\log\det'(\Delta_0),
\end{equation}
которая и приводит к коэффициенту $-\zeta_R(-1)/2$ в одномерном нулевом (по $\mathbb{RP}^3$) секторе KK-ряда по $S^1$.

\paragraph{Шаг 3: вопрос схемозависимости.}
В общем случае ζ-детерминант определяется с масштабом $\mu$:
$\log\det(O/\mu^2)=-\zeta'_O(0)-\zeta_O(0)\log\mu^2$, что потенциально делает константную часть схемозависимой.
На фоне $\mathbb{RP}^3\times S^1$ с метрикой произведения исчезают стандартные 4D локальные инварианты конформной аномалии ($W^2\equiv 0$ и $E_4\equiv 0$), поэтому коэффициент при $\log\mu$ для конформных 4D операторов (Dirac/Maxwell+ghost) обнуляется, а риск схемной неоднозначности существенно уменьшается.
Остаётся технически аккуратно оформить обращение с нулевыми модами (переход к $\det'$ и деление на объём остаточной калибровочной группы).



\paragraph{Типовая схема вывода (что надо сделать).}
Нужно:
\begin{enumerate}
  \item явно выписать спектр соответствующего одномерного оператора на $S^1$ (с учётом калибровочных ограничений и исключения нулевой моды);
  \item построить дзета-функцию (или эквивалентно тепловое ядро) и регуляризовать сумму;
  \item показать, что в калибровочно-инвариантной комбинации (вектор+ghost) остаётся половина одномерного остатка.
\end{enumerate}
Итог должен быть записан как следствие конкретной вычисленной $\zeta_{\mathcal{O}}'(0)$.

\subsection{Механизм поправки второго порядка $\delta_{BB}$}
Термин $\delta_{BB}\propto S_{\mathrm{geo}}^{-2}$ по структуре напоминает вклад следующего порядка по кумулянтам/возмущению.
Чтобы он не выглядел ad hoc, нужно явно указать его происхождение.

\paragraph{Возможные механизмы (что надо выбрать и довести до формулы).}
Например:
\begin{itemize}
  \item вклад второй кумулянты в разложении $\Gamma=-\ln\langle e^{-S_{int}}\rangle$,
  \item смешанные петли (эффекты обратной реакции) и их нормировка через интегралы по $K$,
  \item нелокальный остаток, возникающий после вычитания локальных членов теплового ядра.
\end{itemize}

\paragraph{Кандидатная нормировка, согласованная с геометрией $\mathbb{RP}^3$.}
Если трактовать $\delta_{BB}$ как безразмерный эффект второго порядка (например, как вклад второй кумулянты), то по структуре он должен быть пропорционален квадрату ``малого параметра''.
В нашей постановке естественный малый параметр задаётся отношением единичного квантового вклада к геометрическому ядру $S_{\mathrm{geo}}$, то есть $\sim 1/S_{\mathrm{geo}}$.
Отсюда возникает масштабирование $\delta_{BB}\sim 1/S_{\mathrm{geo}}^2$.

Остаётся геометрическая нормировка безразмерного коэффициента.
При $R=1$ объём $\mathbb{RP}^3$ равен
$\mathrm{Vol}(\mathbb{RP}^3)=\pi^2$, следовательно
\begin{equation}
\mathrm{Vol}(\mathbb{RP}^3)^2=(\pi^2)^2=\pi^4.
\end{equation}
Это даёт естественный кандидат для чисто геометрического подавления второго порядка:
\begin{equation}
\delta_{BB}=\frac{C}{\mathrm{Vol}(\mathbb{RP}^3)^2\,S_{\mathrm{geo}}^2}
=\frac{C}{\pi^4\,S_{\mathrm{geo}}^2},
\end{equation}
где $C$ --- безразмерная константа порядка единицы.
В формуле, используемой в данной работе, принимается наиболее ``жёсткая'' (не вводящая дополнительной ручки) нормировка $C=1$.
Строгий вывод $C=1$ потребует явного 2-loop вычисления в выбранной постановке.


\section{Вариационный принцип для отбора $K$ (программа)}
Для превращения выбора $K$ в следствие теории предлагается рассматривать функционалы спектральной геометрии.

\subsection{Кандидатные функционалы}
Например,
\begin{align}
\mathcal{F}_1[K] &= \zeta_{D}'(0) + \zeta_{\Delta_1}'(0) - 2\,\zeta_{\Delta_0}'(0), \\
\mathcal{F}_2[K] &= \eta_{\mathrm{APS}}(D_K), \\
\mathcal{F}_3[K] &= \mathrm{sys}(K) = \min_{\gamma\in\pi_1(K)\setminus\{e\}} \ell(\gamma).
\end{align}
Здесь комбинация в $\mathcal{F}_1$ отражает типичную калибровочно-инвариантную 1-loop структуру (вектор + ghost).

\subsection{Тестовый набор геометрий}
В качестве первичного теста предлагается сравнить значения $\mathcal{F}[K]$ на семействе
\begin{equation}
K\in\{S^3\times S^1,\ \mathbb{RP}^3\times S^1,\ L(3,1)\times S^1,\ L(5,2)\times S^1\},
\end{equation}
и проверить, достигается ли минимум на $\mathbb{RP}^3\times S^1$ при фиксированных физических ограничениях (spin-структура, нетривиальная $\pi_1$, постоянная положительная кривизна).

\section{Эффективное действие и связь с $\alpha(0)$}
В этой работе принимается рабочая гипотеза: в низкоэнергетическом пределе (IR) обратная константа связи электромагнитного сектора определяется геометрически как значение эффективного действия\footnote{Здесь мы принимаем рабочую гипотезу $\alpha^{-1}(0)\sim\Gamma_{\mathrm{eff}}$. Строгое обоснование этой связи требует восстановления полной структуры КЭД из компактификации и выбора условий ренормировки, что выходит за рамки данной работы.},
\begin{equation}
\alpha^{-1}(0) \sim \Gamma_{\mathrm{eff}}[K] \quad \text{(после фиксации схемы регуляризации и контрчленов)}.
\end{equation}

\subsection{Формула как проверяемое следствие}
При конкретном выборе постановки (см. обсуждение в дедуктивной и строгой ветках проекта) получается выражение
\begin{equation}
\alpha^{-1}
=
\left(4\pi^3+\pi^2+\pi\right)
-
\frac{1}{24\left(4\pi^3+\pi^2+\pi\right)}
-
\frac{1}{\pi^4\left(4\pi^3+\pi^2+\pi\right)^2}.
\end{equation}
На уровне журнальной публикации данная формула должна сопровождаться: (i) явным спектральным вычислением коэффициента $1/24$, (ii) выводом механизма второго порядка, (iii) анализом зависимости от spin-структуры/BC и от схемы вычитания локальных вкладов.

\section{Связь с ренормгруппой и стандартной КТП (сопоставление масштабов)}
Формула в этой работе нацелена на $\alpha(0)$, то есть на предельное низкоэнергетическое значение электромагнитной связи.
Это принципиально: в Стандартной модели $\alpha(\mu)$ ``бежит'' и на характерном масштабе $m_Z$ имеет другое значение.

\paragraph{Почему сравнение идёт с $\alpha(0)$, а не с $\alpha(m_Z)$.}
В нашей постановке вычисляется эффективное действие на компактном внутреннем факторе $K$ и затем фиксируется низкоэнергетическая (IR) нормировка.
Физически это соответствует тому, что при $\mu\to 0$ все массивные поля декуплируются, а остаётся только безмассовый фотон, так что наблюдаемая связь --- это именно $\alpha(0)$.

\paragraph{Стандартное running в КЭД (1-loop).}
На уровне одной петли для КЭД справедливо
\begin{equation}
\beta(\alpha):=\frac{d\alpha}{d\ln\mu}=\frac{2\alpha^2}{3\pi}\sum_f Q_f^2
\end{equation}
(сумма по фермионам, ``активным'' при данном $\mu$).
Эквивалентно, для обратной константы связи
\begin{equation}
\alpha^{-1}(\mu)=\alpha^{-1}(\mu_0)-\frac{2}{3\pi}\sum_f Q_f^2\,\ln\frac{\mu}{\mu_0}+\cdots.
\end{equation}
Переход от $\alpha^{-1}(0)\approx 137.036$ к $\alpha^{-1}(m_Z)\approx 128$ является следствием стандартного ренормгруппового бега с учётом всех заряженных степеней свободы Стандартной модели.

\paragraph{Численная проверка (скрипт проекта).}
Для контроля согласованности ``геометрическая ИК-нормировка $\to$ стандартный РГ-бег'' в проекте используется скрипт \texttt{scripts/running\_alpha.py}.
Он закрывает две разные (и обе необходимые) проверки:
\begin{enumerate}
  \item \textbf{RG-тест коэффициента (спектральный тест).}
  Скрипт печатает (кусочно по порогам) величины
  \begin{equation}
  b(\mu)= -\frac{4}{3}\sum_f Q_f^2,\qquad
  \frac{d\alpha^{-1}}{d\ln\mu}= -\frac{2}{3\pi}\sum_f Q_f^2,
  \end{equation}
  которые должны воспроизводиться в спектральном выводе как коэффициенты при $\ln\mu$ в калибровочно-инвариантной 1-loop комбинации детерминантов (через $\zeta(0)$ или $a_4$).

  \item \textbf{Консистентность геомодели (подбор $R$).}
  Скрипт включает ``геометрическую'' формулу из этой статьи
  \begin{equation}
  \alpha^{-1}_{\mathrm{geo}}(R)=S_{\mathrm{geo}}(R)-\frac{1}{24\,S_{\mathrm{geo}}(R)}-\frac{1}{\pi^4\,S_{\mathrm{geo}}(R)^2},
  \end{equation}
  и позволяет автоматически найти $R_\ast$ из условия $\alpha^{-1}_{\mathrm{geo}}(R_\ast)=\alpha^{-1}_{\mathrm{target}}$.
  Это превращает ``$R\approx1$'' в проверяемую численную подпроверку (пусть пока и не заменяющую строгую вариационную стабилизацию).
\end{enumerate}

\subsection*{Эскиз: спектральный лагранжиан и бегущая константа связи}
Ниже зафиксируем минимальный (пока схематический) вывод того, как $\mu$-зависимость калибровочной константы возникает из спектрального действия и 1-loop детерминантов.
Цель --- иметь в тексте ``мост'' между геометрическим определением ИК-нормировки и стандартной РГ-структурой.

\paragraph{Шаг 0: спектральное действие как эффективный лагранжиан.}
В подходе спектрального действия (Connes--Chamseddine) для 4-мерной евклидовой геометрии вводится функционал
\begin{equation}
S_{\mathrm{SA}}(\Lambda)=\mathrm{Tr}\, f\!\left(\frac{D^2}{\Lambda^2}\right),
\end{equation}
где $D$ --- оператор Дирака (с ``внутренним'' конечным фактором), $\Lambda$ --- УФ-срез, а $f$ --- сглаживающая функция.
Его асимптотическое разложение по тепловому ядру имеет вид
\begin{equation}
\mathrm{Tr}\, f\!\left(\frac{D^2}{\Lambda^2}\right)
\sim
\sum_{n\ge 0} f_{4-n}\,\Lambda^{4-n}\, a_n(D^2),
\end{equation}
где $a_n$ --- коэффициенты теплового ядра, а $f_{k}$ --- моменты функции $f$.

\paragraph{Шаг 1: коэффициент перед $F_{\mu\nu}F^{\mu\nu}$ и определение $g^{-2}(\Lambda)$.}
Калибровочный кинетический член появляется в $a_4$ (в 4D) и имеет общую структуру
\begin{equation}
S_{\mathrm{gauge}}\;\supset\; \frac{1}{4 g^2(\Lambda)}\int d^4x\,\sqrt{g}\; F_{\mu\nu}F^{\mu\nu},
\end{equation}
где $g^{-2}(\Lambda)$ выражается через комбинацию моментов $f$ и спектральных данных внутреннего (конечного) сектора.
В рамках нашей программы это место, где геометрическая нормировка на $K$ задаёт ИК-граничное условие $g^{-2}(\mu_0)$ (а не УФ-значение).

\paragraph{Шаг 2: 1-loop детерминанты и логарифмический бег.}
После интегрирования квантовых флуктуаций возникает стандартная 1-loop добавка
\begin{equation}
\Gamma_{\mathrm{1\text{-}loop}}(\mu)
=\frac{1}{2}\ln\det'\Delta_1-\ln\det'\Delta_0-\ln\det D+\cdots,
\end{equation}
и при введении масштаба $\mu$ через ζ-регуляризацию
\begin{equation}
\ln\det\!\left(\frac{\mathcal{O}}{\mu^2}\right)= -\zeta_{\mathcal{O}}'(0)-\zeta_{\mathcal{O}}(0)\,\ln \mu^2.
\end{equation}
Именно коэффициент при $\ln\mu$ (то есть комбинация $\zeta_{\mathcal{O}}(0)$ для набора операторов Maxwell+ghost+fermions) воспроизводит логарифмический бег $g(\mu)$.

\paragraph{Шаг 3: сопоставление со стандартной формой и вопрос знака.}
В удобной (но конвенциональной) записи
\begin{equation}
\frac{1}{g^2(\mu)}
=
\frac{1}{g^2(\mu_0)}
+
\frac{b}{8\pi^2}\ln\frac{\mu}{\mu_0}
+\cdots,
\end{equation}
где $b$ определяется суммой вкладов полей.
Для КЭД известно, что $\beta(\alpha)>0$ (заряд экранируется), поэтому $\alpha$ растёт с $\mu$, а значит $\alpha^{-1}$ убывает с $\ln\mu$.
В этой конвенции это означает $b<0$ для одного заряженного дираковского фермиона.

\paragraph{Шаг 4: проверка на одном дираковском фермионе (контрольная точка).}
Для одного дираковского фермиона заряда $Q$ стандартная 1-loop формула даёт
\begin{equation}
\alpha^{-1}(\mu)=\alpha^{-1}(\mu_0)-\frac{2}{3\pi}Q^2\,\ln\frac{\mu}{\mu_0}+\cdots.
\end{equation}
Так как $\alpha=g^2/(4\pi)$, получаем эквивалентно
\begin{equation}
\frac{1}{g^2(\mu)}=\frac{1}{g^2(\mu_0)}-\frac{Q^2}{6\pi^2}\,\ln\frac{\mu}{\mu_0}+\cdots,
\end{equation}
то есть в форме выше
\begin{equation}
\frac{b}{8\pi^2}=-\frac{Q^2}{6\pi^2}
\qquad\Longleftrightarrow\qquad
b=-\frac{4}{3}Q^2.
\end{equation}
Эта формула является тем самым ``спектральным тестом'': вычисленная из $\zeta_{\mathcal{O}}(0)$ величина $b$ обязана воспроизводить $-\frac{4}{3}\sum_f Q_f^2$ (с учётом кратностей и порогов).

\paragraph{Шаг 5: как именно это должно получиться спектрально.}
Спектрально коэффициент $b$ извлекается как коэффициент при $\int F_{\mu\nu}F^{\mu\nu}$ в калибровочно-инвариантной комбинации $a_4$ (или эквивалентно $\zeta(0)$) для набора операторов.
В минимальной постановке $U(1)$ это означает: нужно выписать вклад фермионного оператора Дирака (с зарядом $Q$ в ковариантной производной) и вклад Maxwell+ghost, затем собрать их в одну комбинацию и сравнить результат с контрольной точкой выше.

\paragraph{Что в этой статье уже сделано, а что остаётся.}
В тексте уже фиксируется кандидатная ИК-нормировка через $\Gamma_{\mathrm{eff}}[K]$ и обсуждается происхождение одномерного остатка $1/24$.
Однако для настоящего ``спектрального вывода бега'' пока не хватает явного вычисления коэффициента при $\ln\mu$ в 1-loop эффективном действии.

\paragraph{Технический недостающий элемент: $\zeta_{\mathcal{O}}(0)$ и коэффициенты теплового ядра.}
Для лапласовского оператора на 4-мерном фоне (в частности, для калибровочно-инвариантной комбинации Maxwell+ghost и для фермионного оператора) ζ-регуляризация даёт
\begin{equation}
\ln\det\!\left(\frac{\mathcal{O}}{\mu^2}\right)= -\zeta_{\mathcal{O}}'(0)-\zeta_{\mathcal{O}}(0)\,\ln \mu^2,
\end{equation}
и именно $\zeta_{\mathcal{O}}(0)$ задаёт логарифмическую $\mu$-зависимость.
При этом в 4D (для операторов типа Лапласа) имеется связь
\begin{equation}
\zeta_{\mathcal{O}}(0)=a_4(\mathcal{O})-\dim\ker(\mathcal{O}),
\end{equation}
где $a_4(\mathcal{O})$ --- коэффициент теплового ядра (локальный инвариант).
Таким образом, задача ``вывести running'' сводится к вычислению нужной калибровочно-инвариантной комбинации $a_4$ на выбранном фоне и к корректному учёту нулевых мод (то есть переходу к $\det'$).

\paragraph{Минимальный план вычисления коэффициента $b$ (без расширения модели).}
\begin{enumerate}
  \item \textbf{Зафиксировать фон и схему:} 4D фон (плоский $M_4$), внутренняя геометрия $K=\mathbb{RP}^3\times S^1$ с заданной метрикой произведения, а также правило вычитания локальных контрчленов.

  \item \textbf{Выбрать набор операторов:}
  \begin{equation}
  \Gamma^{(1)}=\tfrac12\ln\det'\Delta_1-\ln\det'\Delta_0-\ln\det D+\cdots
  \end{equation}
  (и уточнить, какие фермионные степени свободы считаются ``активными'' при данном $\mu$).

  \item \textbf{Вычислить $\zeta(0)$ через $a_4$:} извлечь коэффициент при $F_{\mu\nu}F^{\mu\nu}$ из суммарного $a_4$ для выбранной комбинации операторов.

  \item \textbf{Пороговое согласование:} реализовать декуплирование массивных фермионов как кусочно-логарифмическую зависимость (меняется набор ``активных'' $Q_f$).

  \item \textbf{Сопоставить со стандартной формой:}
  \begin{equation}
  \frac{1}{g^2(\mu)}
  =
  \frac{1}{g^2(\mu_0)}
  +
  \frac{b}{8\pi^2}\ln\frac{\mu}{\mu_0}
  +\cdots,
  \end{equation}
  и проверить, что $b$ совпадает с 1-loop значением, эквивалентным формуле для $\beta(\alpha)$ в КЭД.
\end{enumerate}

\paragraph{Что остаётся сделать на уровне строгого сопоставления.}
Для журнального статуса необходимо:
\begin{itemize}
  \item чётко зафиксировать условие ренормировки, в котором величина $\Gamma_{\mathrm{eff}}[K]$ интерпретируется как $\alpha^{-1}(0)$;
  \item проверить, что при переходе к $\alpha(\mu)$ при $\mu>m_e$ стандартное running воспроизводится и не требует модификации КЭД;
  \item описать роль других секторов (не-$U(1)$) и объяснить, почему топология $\pi_1(K)$ на минимальном уровне влияет именно на электромагнитную нормировку.
\end{itemize}

\section{Независимые предсказания (конъектуры и тесты)}
Чтобы отделить ``совпадение с $\alpha$'' от подгонки, модель должна давать проверяемые следствия, не использующие $\alpha$ как вход.
Как возможные направления тестов фиксируем:
\begin{itemize}
  \item спектральные отношения (с кратностями) для операторов на $\mathbb{RP}^3$ как источник дискретных иерархий;
  \item геометрические оценки углов смешивания как функций отношения ``фермионного'' и ``бозонного'' вкладов (с обязательной проверкой устойчивости к выбору BC/схемы);
  \item ограничения на топологические фазы/параметры вакуума при $\pi_1(\mathbb{RP}^3)=\mathbb{Z}_2$.
\end{itemize}
Эти пункты в текущей версии формулируются как программа, поскольку требуют явной привязки к наблюдаемым (массы, углы, ЭДМ) и контроля систематик.

\section{Критические вопросы (дорожная карта)}
Ниже собран список вопросов, которые фиксируют ``белые пятна'' между численным совпадением и строгим выводом, в порядке от математической самосогласованности к физической интерпретации и феноменологии.

\subsection*{Начало закрытия вопросов (частичные ответы)}
Ниже мы фиксируем минимальные ответы/наброски ответов, которые уже можно оформить в статье без выхода за рамки честного статуса ``программы''.
Цель этого блока --- заменить часть вопросов на проверяемые утверждения (с явными местами, где требуется расчёт).

\paragraph{(1) Почему спинорный сектор ``видит'' накрытие $S^3$? (уточнение механизма).}
Спин-структура на $\mathbb{RP}^3=S^3/\mathbb{Z}_2$ удобно описывается через подъём на накрытие: спинорные поля на $\mathbb{RP}^3$ эквивалентны $\mathbb{Z}_2$-эквивариантным спинорам на $S^3$ (с проекцией на чётный/нечётный сектор относительно антиподального отображения).
В функциональном интеграле это означает, что трасса/детерминант в спинорном секторе может быть записан как
\emph{трасса на накрытии с вставкой проектора} на нужный $\mathbb{Z}_2$-сектор.
В такой форме появляются две конкурирующие ``половины'': (i) фактор $1/|\mathbb{Z}_2|$ от уменьшения фундаментальной области и (ii) фактор 2 от удвоения спектральных данных при переходе к оператору на накрытии с проектором.
В результате нормировка, выраженная через интеграл по плотности состояний, естественным образом приводит к масштабу, пропорциональному $\mathrm{Vol}(S^3\times S^1)$.
Строго это должно быть оформлено как лемма о связи ζ-детерминантов на $S^3$ и на $\mathbb{RP}^3$ через проекцию по $\mathbb{Z}_2$.

\paragraph{(2) Коэффициент $1/24$: что именно остаётся проверить.}
В тексте уже дан кандидатный вывод $1/24$ как одномерного ζ-регуляризованного остатка и объяснено происхождение множителя $1/2$ из калибровочно-инвариантной комбинации Maxwell+ghost.
Для закрытия вопроса остаётся выполнить (и явно прописать) два технических шага:
(а) аккуратно обосновать факторизацию (или выделение 1D-нулевого сектора) в ζ-детерминанте на произведении $\mathbb{RP}^3\times S^1$;
(б) показать, что кривизна $\mathbb{RP}^3$ не изменяет конечную константную часть именно в рассматриваемой комбинации (в частности, что возможные локальные контрчлены либо отсутствуют, либо фиксируются выбранным принципом ренормировки).

\paragraph{(3) Вариационная устойчивость $L(2,1)$.}
``Быстрый отбор'' в классе сферических форм показывает, что при фиксации ключевых нормировок (объёмный критерий $\mathrm{Vol}(M^3)=\pi^2$ и дискретность $U(1)$-flat без непрерывных модулей) вариант $M^3=\mathbb{RP}^3=L(2,1)$ выделяется уже на уровне грубой топологии.
Чтобы закрыть вопрос о глобальном минимуме, требуется вычислить выбранный функционал (например, $\mathcal{F}_1[K]=\zeta_{D}'(0)+\zeta_{\Delta_1}'(0)-2\zeta_{\Delta_0}'(0)$) для семейства $L(p,q)\times S^1$ и сравнить значения при одинаковых физических ограничениях (spin-структура, BC по $S^1$, обращение с нулевыми модами).
До выполнения этого расчёта корректный статус остаётся таким: $\mathbb{RP}^3$ --- наиболее жёстко фиксируемый кандидат, но не доказанный глобальный минимум.

\subsection{Блок 1: геометрическая самосогласованность}
\begin{enumerate}
  \item \textbf{Обоснование смешивания объёмов накрытия и фактора.}
  В $S_{\mathrm{geo}}$ используются разные геометрические объёмы: $4\pi^3=\mathrm{Vol}(S^3\times S^1)$ для фермионов и $\pi^2=\mathrm{Vol}(\mathbb{RP}^3)$ (с последующей нормировкой по $S^1$) для бозонного сектора.
  \emph{Вопрос:} каков строгий физический механизм, заставляющий фермионы ``чувствовать'' объём накрытия $S^3$, а не физического фактора $\mathbb{RP}^3$? Почему принцип локальности не ограничивает интегрирование фундаментальной областью фактора и не превращает этот шаг в скрытую подгонку?

  \item \textbf{Выбор спектральной задачи для коэффициента $1/24$.}
  Коэффициент $1/24$ мотивируется одномерным касимировским (или ζ-регуляризованным) остатком.
  \emph{Вопрос:} выполнен ли прямой расчёт ζ-детерминантов для соответствующих операторов (Maxwell+ghost, Dirac) именно на $\mathbb{RP}^3\times S^1$ с выбранной spin-структурой и условиями по $S^1$? Учтены ли особенности фактора $\mathbb{RP}^3$ (включая роль торсиона в когомологиях, обращение с нулевыми модами и возможную чувствительность к схеме вычитания локальных членов теплового ядра), и не меняет ли кривизна 3D-фактора конечную константную часть по сравнению с плоским цилиндром?

  \item \textbf{Вариационная устойчивость $L(2,1)$.}
  Геометрия $\mathbb{RP}^3=L(2,1)$ заявляется как выделенная.
  \emph{Вопрос:} сравнивалось ли значение выбранного функционала (или эффективного действия) $\mathcal{F}[K]$ для $K=L(p,q)\times S^1$ при различных $(p,q)$? Если, например, для $L(3,1)$ или $L(5,2)$ вакуумная энергия ниже, то возникает вопрос о туннелировании и изменении предсказываемого значения $\alpha$; есть ли аргумент, что $L(2,1)$ даёт глобальный минимум в физически допустимом классе?
\end{enumerate}

\subsection*{Продолжение закрытия вопросов: блок 2 (частичные ответы)}

\paragraph{(4) Две spin-структуры на $\mathbb{RP}^3$ и ``три поколения''.}
Факт: на $\mathbb{RP}^3$ существует две spin-структуры, то есть в чисто геометрической постановке есть лишь два дискретных варианта спектра Дирака.
На текущем этапе это следует трактовать не как готовое объяснение трёх поколений, а как \emph{ограничение на минимальную модель}: ``топологический индекс'' поколения не может быть просто выбором spin-структуры на одном и том же $\mathbb{RP}^3$.
Естественные (и проверяемые) варианты закрытия этого разрыва:
\begin{itemize}
  \item поколения порождаются не геометрией $\mathbb{RP}^3$, а внутренним конечным спектральным пространством (дискретными матричными алгебрами НКГ), тогда spin-структура отвечает лишь за ``класс'' фермионного сектора (например, за наличие/отсутствие определённых мод);
  \item три поколения реализуются как три сектора граничных условий/голономий по $S^1$ (например, разные допустимые фазы для спиноров), при этом spin-структура $\mathbb{RP}^3$ фиксируется однажды;
  \item либо геометрия должна быть расширена (например, $L(p,q)$ с большим $p$), и тогда вопрос о выборе $\mathbb{RP}^3$ должен решаться вариационным принципом (см. Блок 1).
\end{itemize}
В любом из вариантов здесь требуется не обсуждение, а явная спектральная проверка: как именно дискретные данные (spin-структура/голономии/конечная алгебра) влияют на кратности мод и на нормировку 4D-теории.

\paragraph{(5) Стабилизация $R=1$: какой минимум вообще должен быть доказан.}
Чтобы закрыть вопрос, нужно заменить утверждение ``берём $R=1$'' на строгое:
\begin{equation}
\left.\partial_R V_{\mathrm{eff}}(R)\right|_{R=R_\ast}=0,
\qquad
\left.\partial_R^2 V_{\mathrm{eff}}(R)\right|_{R=R_\ast}>0,
\end{equation}
где $R_\ast$ затем сопоставляется с единицей в выбранной системе единиц.
В рамках спектрального действия/NКГ естественным кандидатом является функционал типа
$S_{\mathrm{SA}}\sim\mathrm{Tr}\,f(D^2/\Lambda^2)$ (плюс 1-loop детерминанты), дающий вклад $\propto\mathrm{Vol}(K)\Lambda^4$ и кривизновые поправки.
Даже без полного расчёта минимально необходимый шаг ``закрытия'' здесь: выписать зависимость ведущих членов (объёмный, скалярная кривизна, Касимир) от $R$ и проверить, что конкуренция степеней $R$ допускает локальный минимум.

\paragraph{(6) Почему сравнение идёт с $\alpha(0)$, если геометрия планковская.}
В этой работе геометрическое число интерпретируется как \emph{граничное условие на ИК-нормировку} (renormalization condition) для $U(1)$-сектора: вычисляется $\alpha(0)$, после чего стандартный ренормгрупповой бег переносит значение к любому $\mu$.
Тем самым ``планковость'' геометрии не означает, что сравнение делается с УФ-константой: геометрия выступает как микроскопический механизм фиксации константы, а измеряемая величина задаётся выбором схемы ренормировки.
Чтобы закрыть вопрос строго, нужно (i) явно сформулировать условие ренормировки, связывающее $\Gamma_{\mathrm{eff}}[K]$ с $\alpha^{-1}(0)$, и (ii) показать, что добавление стандартного running не требует новых свободных параметров.

\subsection{Блок 2: физическая интерпретация и Стандартная модель}
\begin{enumerate}
  \setcounter{enumi}{3}
  \item \textbf{Проблема поколений фермионов (spin-структуры).}
  Для $\mathbb{RP}^3$ существует две неэквивалентные spin-структуры.
  \emph{Вопрос:} как это согласуется с тремя поколениями фермионов в Стандартной модели? Предполагается ли, что одно поколение имеет иную природу (составность/дополнительный внутренний индекс/другой топологический сектор), или же сама геометрия должна быть модифицирована (например, замена на $L(3,1)$)? Как именно распределяются спинорные поля по spin-структурам и как это отражается в спектре и в 4D-теории?

  \item \textbf{Стабилизация модуля $R=1$.}
  В тексте используется $R=1$ (в планковских единицах), но для предсказательного статуса нужен динамический механизм.
  \emph{Вопрос:} какой вклад в $V_{\mathrm{eff}}(R)$ (Казимир, потоки, гравитационные/квантовые поправки, дополнительные поля) обеспечивает устойчивый минимум именно при $R=1$? Почему в данной постановке радиус не уходит в $0$ или $\infty$ без дополнительной структуры (например, SUSY или стабилизации модулей)?

  \item \textbf{Иерархия масштабов и бегущая константа.}
  Геометрия на $R\sim \ell_{\mathrm{Pl}}$ естественно относится к УФ, тогда как сравнение делается с $\alpha(0)$.
  \emph{Вопрос:} почему геометрический вклад, сформированный на планковском масштабе, не ``размывается'' при ренормгрупповом спуске на много порядков энергии, и как именно учитывается (или обосновывается пренебрежение) вклад $\ln(M_{\mathrm{Pl}}/m_e)$ при переходе от УФ-нормировки к ИК-наблюдаемой $\alpha(0)$?
\end{enumerate}

\subsection*{Продолжение закрытия вопросов: блок 3 (частичные ответы)}

\paragraph{(7) Остальные взаимодействия: почему здесь пока речь только об $U(1)$.}
Текущее совпадение сконструировано как ИК-нормировка именно электромагнитного сектора и использует топологические особенности $\pi_1(K)$ (дискретные голономии) и калибровочно-инвариантную 1-loop комбинацию Maxwell+ghost.
В таком виде конструкция по определению более ``чувствительна'' к абелеву сектору, где голономии напрямую параметризуют физически различимые плоские связи.
Чтобы закрыть вопрос, нужно расширить постановку до полноценного спектрального действия с конечным пространством НКГ (описывающим $SU(2)\times SU(3)$) и показать, что:
(i) те же принципы нормировки фиксируют не только $\alpha$, но и отношения констант связей на некотором масштабе;
(ii) при РГ-беге получается стандартная картина.
До этого корректный статус: модель пока претендует на выделение $\alpha(0)$ как тестового ``маячка'', а не на полный вывод параметров Стандартной модели.

\paragraph{(8) Монополи и топологические дефекты: что можно сказать уже сейчас.}
Сам по себе фактор $\mathbb{RP}^3$ не вводит новые непрерывные модульные параметры для $U(1)$-плоских связей (в отличие от многообразий с $b_1>0$), что уменьшает риск появления нежелательных лёгких мод.
В то же время наличие дискретных голономий в принципе допускает существование топологических секторов.
Минимальный шаг ``закрытия'' здесь --- явно выписать классификацию $U(1)$-связей и топологических зарядов на $K$ (через соответствующие (ко)гомологические группы) и оценить характерный масштаб масс дефектов как $\sim 1/R$.
Если $R\sim \ell_{\mathrm{Pl}}$, то любые такие объекты оказываются планковски тяжёлыми и не противоречат текущим экспериментальным ограничениям; если же стабилизация даёт иной масштаб, этот вывод должен быть пересмотрен.

\paragraph{(9) ``Топологическая защита протона'': что требуется для утверждения.}
Топология $\pi_1(\mathbb{RP}^3)=\mathbb{Z}_2$ сама по себе не запрещает барион-нарушающие операторы в 4D: запрет должен следовать либо из симметрий конечного НКГ-сектора (алгебра наблюдаемых и её внутренние автоморфизмы), либо из селекционных правил, связанных с дискретными фазами/голономиями.
Минимально проверяемое утверждение, которое можно поставить как задачу: вывести, появляются ли в эффективном 4D действии операторы размерности 6 типа $qqql/\Lambda^2$ при интегрировании по $K$, и если да --- какие коэффициенты и какие селекционные правила диктует дискретная топология.

\paragraph{Как использовать феноменологический эскиз ``эволюции вакуума'' (что оставить, что вычеркнуть).}
Ниже фиксируем ``фильтр'' для расширенного (пока спекулятивного) текста, чтобы он работал как источник зацепок, но не превращал статью в набор неподтверждённых утверждений.

\paragraph{Можно сохранить как совместимые зацепки (в текущем статусе работы).}
\begin{itemize}
  \item \textbf{Аксиоматическую цепочку топологии} $\mathbb{H}\Rightarrow S^3\Rightarrow \mathbb{RP}^3$ (фактор по центральному $\mathbb{Z}_2$) и \textbf{периодичность времени} $T\cong S^1$ как мотивацию выбора $K=\mathbb{RP}^3\times S^1$.
  \item \textbf{Разделение на ``геометрическое ядро + квантовую поправку''} как удобный язык для организации вычислений (при условии, что каждое слагаемое привязано к конкретному оператору/детерминанту и схеме вычитания).
  \item \textbf{Идею трёх режимов (IR/электрослабый/GUT)} как качественную карту того, что надо сопоставлять через стандартный running и возможную смену эффективного описания.
\end{itemize}

\paragraph{Нужно вычеркнуть либо жёстко переформулировать как гипотезы (иначе это будет подгонка).}
\begin{itemize}
  \item Любые формулировки вида ``мы доказали'' / ``Q.E.D.'': их следует заменить на ``предлагаем модель'' / ``гипотеза''.
  \item Введение подавляющего фактора $\bigl(1+\xi\ln^6(\mu/m_e)\bigr)^{-1}$: параметр $\xi$ и степень 6 должны быть \emph{выведены} (например, из спектрального действия/конкретного 1-loop/2-loop расчёта или из модулярной динамики), иначе это дополнительная ручка.
  \item ``Плавление размерности $4\to 2$'' и привязка к ``компактификации на Калаби--Яу'' без явного перехода в геометрии/спектре: это допустимо только как \emph{сценарий} с перечислением проверяемых критериев (например, изменение спектральной размерности по тепловому ядру).
  \item Абсолютная величина и квадратичные ``струнные'' поправки в $\Delta_{\mathrm{quant}}$: нужен знак, происхождение (какие диаграммы/детерминанты) и контроль размерностей.
  \item Численные утверждения про $\alpha^{-1}(m_Z)$ и тем более про GUT ($\sim 25$) как ``полное совпадение'': без прозрачного вычисления (и без явной зависимости от выбранных допущений) это следует оставить только как ориентиры, а не как вывод.
\end{itemize}

\paragraph{Шаг вперёд (как превратить эскиз в часть строгой статьи).}
Минимальный вариант: заменить все новые ``законы эволюции'' на одну чётко определённую величину (например, $\Gamma_{\mathrm{eff}}(R,\mu)$ или спектральное действие $S_{\mathrm{SA}}(R,\Lambda)$), а затем вывести (а не постулировать) (i) $R$-динамику/стабилизацию и (ii) связь с условием ренормировки для $\alpha(0)$.

\subsection{Блок 3: феноменология и проверяемость}
\begin{enumerate}
  \setcounter{enumi}{6}
  \item \textbf{Судьба остальных взаимодействий.}
  Если электромагнетизм фиксируется геометрией через $U(1)$ и структуру $S^1$, остаётся вопрос об остальных калибровочных секторах.
  \emph{Вопрос:} можно ли в рамках той же схемы получить ограничения/значения для $SU(2)$ и $SU(3)$ (например, угол Вайнберга или $\alpha_s$) без введения новых ``ручек''? Если метод работает только для $\alpha$ и не переносится на другие константы, не является ли совпадение численной случайностью?

  \item \textbf{Масса магнитного монополя.}
  Компактная $U(1)$-теория в общем случае допускает монопольные конфигурации.
  \emph{Вопрос:} какой порядок величины массы монополя предсказывает данная компактификация (ожидаемо $\sim 1/R$)? Не приводит ли топология $\mathbb{RP}^3$ и дискретные голономии к появлению более лёгких монополей или иных топологических дефектов, которые уже исключены экспериментально?

  \item \textbf{Топологическая защита протона.}
  \emph{Вопрос:} как инварианты $\pi_1(\mathbb{RP}^3)$ и $\pi_3(\mathbb{RP}^3)$ (а также торсионные классы) отражаются на возможных барион-нарушающих операторах и каналах распада протона? Возникают ли дополнительные ограничения или, наоборот, новые каналы, проверяемые в экспериментах типа Super-Kamiokande?
\end{enumerate}

\section{Заключение}
Переформулирована постановка задачи: из ``формулы, совпадающей с числом'' она переведена в форму проверяемой гипотезы с дорожной картой доказательств.
Следующий необходимый шаг для журнального статуса --- построение $V_{\mathrm{eff}}(R)$ и вариационного критерия для отбора $K$, а также полная спектральная проверка коэффициентов в регуляризованных детерминантах.
Без выполнения хотя бы одного из ключевых пунктов программы (стабилизация $R$, вариационный отбор $K$, спектральный вывод коэффициента $1/24$) гипотеза остаётся эстетически привлекательной, но физически непроверяемой.

\appendix

\section{Ренормусловие для \texorpdfstring{$\alpha(0)$}{alpha(0)} и нормировки}
В этом приложении фиксируем, что именно считается наблюдаемым входом (\texorpdfstring{$\alpha(0)$}{alpha(0)}) и как оно связано с коэффициентом при $\int_{M_4}F_{\mu\nu}F^{\mu\nu}$ в 4D эффективном действии.

\subsection{Нормировка 4D калибровочного члена и определение \texorpdfstring{$\alpha$}{alpha}}
В евклидовой 4-мерной эффективной теории калибровочный кинетический член пишем в виде
\begin{equation}
\Gamma_{\mathrm{eff}}\;\supset\;\frac{1}{4e^2(\mu)}\int_{M_4} d^4x\, F_{\mu\nu}F^{\mu\nu},
\end{equation}
где $e(\mu)$ --- перенормированная константа связи в выбранной схеме при масштабе $\mu$.
Наблюдаемая (в гауссовых единицах) постоянная тонкой структуры определяется как
\begin{equation}
\alpha(\mu):=\frac{e^2(\mu)}{4\pi}.
\end{equation}

\subsection{Томсоновский предел и связь с вакуумной поляризацией}
Пусть 1PI двухточечная функция фотона имеет стандартную структуру
\begin{equation}
\Gamma^{(2)}_{\mu\nu}(q)
=
\bigl(q^2\delta_{\mu\nu}-q_\mu q_\nu\bigr)\,\Pi_{\mathrm{R}}(q^2;\mu)
+\frac{1}{e^2(\mu)}\bigl(q^2\delta_{\mu\nu}-q_\mu q_\nu\bigr).
\end{equation}
Тогда определение ИК-нормировки (Томсоновский предел) формулируется как условие на полюс при $q^2\to 0$:
\begin{equation}\label{eq:renorm_condition_alpha0}
\frac{1}{e^2(0)}
:=
\lim_{q^2\to 0}
\left(\frac{1}{e^2(\mu)}+\Pi_{\mathrm{R}}(q^2;\mu)\right).
\end{equation}
Эквивалентно можно записать это через $\Pi_{\mathrm{R}}(0;\mu)$:
\begin{equation}\label{eq:e0_from_Pi}
\frac{1}{e^2(0)}
=
\frac{1}{e^2(\mu)}+\Pi_{\mathrm{R}}(0;\mu).
\end{equation}
Эти равенства следует понимать как \emph{определение} $e(0)$ в фиксированной схеме; они не предполагают, что $\Pi_{\mathrm{R}}(0;\mu)$ вычислено внутри данной работы.

\subsection{Геометрическая поправка \texorpdfstring{$\Delta_K$}{DeltaK} как вклад внутренних детерминантов}
В рамках программы статьи предполагается, что после интегрирования по внутреннему фактору $K$ и после вычитания локальных 4D контрчленов (контролируемых коэффициентами теплового ядра/значениями $\zeta_{\mathcal{O}}(0)$) остаётся конечная величина $\Delta_K(0;R)$, такая что
\begin{equation}
\frac{1}{e^2(0)}
=
\frac{1}{e^2_{\mathrm{4D}}(0)}
+
\Delta_K(0;R).
\end{equation}
Тем самым, строгость сопоставления с наблюдаемой \texorpdfstring{$\alpha(0)$}{alpha(0)} сводится к (i) вычислению $\Delta_K$ в выбранной постановке (операторы, BC, $\det'$, гармонические моды) и (ii) доказательству, что выбранная схема вычитания действительно соответствует определению $\Pi_{\mathrm{R}}(0;\mu)$ в \eqref{eq:renorm_condition_alpha0}--\eqref{eq:e0_from_Pi}.

\subsection{Метод фонового поля: извлечение поправки к коэффициенту при \texorpdfstring{$F^2$}{F2}}
Чтобы связать 1-loop детерминанты на $M_4\times K$ с поправкой к $e^{-2}(0)$, следует использовать метод фонового поля.
Рассмотрим разложение $A_M=\bar A_M+a_M$, где $\bar A_M$ --- внешний (фоновой) потенциал, а $a_M$ --- квантовая флуктуация.
Мы выбираем фон, который имеет ненулевую 4D компоненту $\bar A_\mu(x)$ и не зависит от внутренних координат, а внутренние компоненты полагаем нулевыми:
\begin{equation}
  \bar A_M=(\bar A_\mu(x),0),
  \qquad
  \partial_y \bar A_\mu=0.
\end{equation}
Тогда 1-loop эффективное действие имеет разложение по инвариантам фонового поля вида
\begin{equation}
  \Gamma^{(1)}[\bar A]
  =
  \Gamma^{(1)}[0]
  +
  \frac{1}{4}\,\Delta\!\left(\frac{1}{e^2}\right)
  \int_{M_4} d^4x\, \bar F_{\mu\nu}\bar F^{\mu\nu}
  +\cdots.
\end{equation}
Определение геометрической поправки $\Delta_K$ в этой постановке состоит в идентификации
\begin{equation}
  \Delta_K(0;R)
  \equiv
  \Delta\!\left(\frac{1}{e^2(0)}\right)
\end{equation}
после вычитания локальных контрчленов (то есть после удаления локальных вкладов, аналитичных по кривизнам и по $\bar F$).

Практически это означает: нужно вычислить коэффициент при $\int_{M_4}\bar F^2$ в калибровочно-инвариантной комбинации
\begin{equation}
  \Gamma^{(1)}_{\mathrm{gauge}}[\bar A]
  =
  \tfrac12\ln\det'\Delta_1(\bar A)
  -\ln\det'\Delta_0(\bar A)
  \;(+\;\text{при необходимости вклад фермионов}),
\end{equation}
и затем перейти к Томсоновской нормировке, то есть к $q^2\to 0$ в смысле \eqref{eq:renorm_condition_alpha0}.

\subsection{Контрольный тест: воспроизведение стандартного коэффициента бега}
В любой схеме, согласованной с КЭД, коэффициент при $\ln\mu$ в 1-loop части обязан воспроизводить известный 1-loop бег.
В ζ-регуляризации
\begin{equation}
  \ln\det\!\left(\frac{\mathcal{O}}{\mu^2}\right)= -\zeta_{\mathcal{O}}'(0)-\zeta_{\mathcal{O}}(0)\,\ln \mu^2,
\end{equation}
поэтому проверка нормировок сводится к вычислению комбинации $\zeta_{\mathcal{O}}(0)$ для набора операторов в присутствии фонового $\bar F$.

Минимальная контрольная точка: один дираковский фермион заряда $Q$ в 4D даёт
\begin{equation}
  \alpha^{-1}(\mu)=\alpha^{-1}(\mu_0)-\frac{2}{3\pi}Q^2\,\ln\frac{\mu}{\mu_0}+\cdots,
\end{equation}
или, эквивалентно,
\begin{equation}
  \frac{1}{e^2(\mu)}=\frac{1}{e^2(\mu_0)}-\frac{Q^2}{6\pi^2}\,\ln\frac{\mu}{\mu_0}+\cdots.
\end{equation}
Следовательно, в нашей внутренней постановке (после редукции на $K$ и вычитания локальных контрчленов) коэффициент при $\ln\mu$ в $\Delta_K$ обязан совпасть с $-Q^2/(6\pi^2)$, с учётом кратностей и порогов.
Если этот тест не выполнен, то сопоставление с наблюдаемой \texorpdfstring{$\alpha(0)$}{alpha(0)} является схемным и требует пересмотра нормировок.

\section{Геометрические леммы для $S^3$, $\mathbb{RP}^3$ и $S^1$}
В этом приложении фиксируем геометрические формулы, которые используются как ``нормировочные'' числа в основной части.

\subsection{Объёмы}
Пусть $S^3$ имеет радиус $R$ (как вложенная гиперсфера в $\mathbb{R}^4$). В гиперсферических координатах метрика имеет вид
\begin{equation}
  ds^2=R^2\bigl(d\chi^2+\sin^2\chi\,d\Omega_2^2\bigr),\qquad \chi\in[0,\pi].
\end{equation}
Тогда
\begin{equation}
  \mathrm{Vol}(S^3)=2\pi^2R^3.
\end{equation}
Так как $\mathbb{RP}^3\cong S^3/\mathbb{Z}_2$ с свободным антиподным действием $x\mapsto -x$, покрытие имеет степень 2 и является локальной изометрией, поэтому
\begin{equation}
  \mathrm{Vol}(\mathbb{RP}^3)=\frac12\,\mathrm{Vol}(S^3)=\pi^2R^3.
\end{equation}

Пусть окружность $S^1$ имеет радиус $R$, тогда её длина
\begin{equation}
  L_{S^1}=\mathrm{Vol}(S^1)=2\pi R.
\end{equation}
Для произведения (в метрике произведения) получаем
\begin{equation}
  \mathrm{Vol}(S^3\times S^1)=\mathrm{Vol}(S^3)\,\mathrm{Vol}(S^1)=4\pi^3R^4.
\end{equation}

\subsection{Систола $\mathbb{RP}^3$}
Нетривиальный класс в $\pi_1(\mathbb{RP}^3)=\mathbb{Z}_2$ соответствует пути на $S^3$, соединяющему точку с антиподом.
Кратчайший такой путь --- геодезическая (дуга большого круга) длины $\pi R$; её проекция в $\mathbb{RP}^3$ даёт кратчайшую нестягиваемую петлю.
Следовательно,
\begin{equation}
  L_{\mathrm{sys}}(\mathbb{RP}^3)=\pi R.
\end{equation}

\section{1-loop структура Maxwell+ghost и роль $\det'$ на $M_4\times K$}
В этом приложении приводим компактный, самодостаточный вывод той части 1-loop архитектуры, которая нужна для однозначного определения комбинаций детерминантов (включая нулевые и гармонические моды).

\subsection{Евклидово действие, калибровочная фиксация и ghost}
Рассматриваем евклидову абелеву теорию на произведении $M_4\times K$:
\begin{equation}
  S_A=\frac{1}{4g_5^2}\int_{M_4} d^4x\int_K d\mathrm{vol}\, F_{MN}F^{MN}.
\end{equation}
Фиксируем ковариантную калибровку $G(A)=\nabla^M A_M=0$ добавлением
\begin{equation}
  S_{\mathrm{gf}}=\frac{1}{2\xi}\int_{M_4} d^4x\int_K d\mathrm{vol}\,(\nabla^M A_M)^2.
\end{equation}
Для $U(1)$ имеем $A^\omega=A+\nabla\omega$, поэтому оператор Фаддеева--Попова равен $\nabla^2$ и возникает ghost-сектор со скалярным лапласианом
\begin{equation}
  S_{\mathrm{gh}}=\int_{M_4} d^4x\int_K d\mathrm{vol}\,\bar c\,\Delta_{\mathrm{gh}}\,c,
  \qquad \Delta_{\mathrm{gh}}=-\nabla^2.
\end{equation}

\subsection{Гауссовы интегралы и знаки детерминантов}
Квадратичная часть по флуктуациям даёт стандартные гауссовы интегралы:
\begin{equation}
  \int\mathcal{D}\varphi\,\exp\bigl(-\tfrac12\langle\varphi,O\varphi\rangle\bigr)\propto (\det O)^{-1/2}
  \quad\Rightarrow\quad
  \Gamma\supset +\tfrac12\ln\det O,
\end{equation}
а для ghost (Grassmann)
\begin{equation}
  \int\mathcal{D}\bar c\,\mathcal{D}c\,\exp\bigl(-\langle\bar c,\Delta_{\mathrm{gh}}c\rangle\bigr)\propto \det\Delta_{\mathrm{gh}}
  \quad\Rightarrow\quad
  \Gamma\supset -\ln\det\Delta_{\mathrm{gh}}.
\end{equation}
Если присутствует заряжённый дираковский фермион $\Psi$,
\begin{equation}
  \int\mathcal{D}\bar\Psi\,\mathcal{D}\Psi\,\exp\bigl(-\langle\bar\Psi,D\!\!\!/\,\Psi\rangle\bigr)\propto \det(D\!\!\!/\,)
  \quad\Rightarrow\quad
  \Gamma\supset -\ln\det(D\!\!\!/\,).
\end{equation}
Итого,
\begin{equation}
  \Gamma^{(1)}=\tfrac12\ln\det(\Delta_1)-\ln\det(\Delta_0)-\ln\det(D\!\!\!/\,)+\cdots,
\end{equation}
где $\Delta_1$ --- оператор на 1-формах (калибровочном поле), а $\Delta_0$ --- скалярный лапласиан (ghost).

\subsection{Разложение Ходжа и калибровочно-инвариантная комбинация}
На компактном многообразии разложение Ходжа для 1-форм имеет вид
\begin{equation}
  A=A_{\perp}+d\phi + A_{\mathrm{harm}},
  \qquad d^{\ast}A_{\perp}=0,
  \qquad A_{\mathrm{harm}}\in\ker \Delta_1.
\end{equation}
На продольных модах $d\phi$ действие $\Delta_1$ редуцируется к скалярному лапласиану, поэтому на уровне спектра (с оговорками о нулевых модах) выполняется факторизация
\begin{equation}
  \det'(\Delta_1)=\det'(\Delta_{1,\perp})\,\det'(\Delta_0).
\end{equation}
Отсюда калибровочно-инвариантная 1-loop комбинация Maxwell+ghost сводится к
\begin{equation}
  \tfrac12\ln\det'(\Delta_1)-\ln\det'(\Delta_0)
  =\tfrac12\ln\det'(\Delta_{1,\perp})-\tfrac12\ln\det'(\Delta_0),
\end{equation}
то есть ghost-сектор исключает продольные степени свободы.

\subsection{Нулевые и гармонические моды на $K=\mathbb{RP}^3\times S^1$}
Для $K=\mathbb{RP}^3\times S^1$ имеем (над $\mathbb{R}$)
\begin{equation}
  b_0(K)=1,\qquad b_1(K)=1.
\end{equation}
Это означает:
\begin{itemize}
  \item в ghost-секторе $\Delta_0$ всегда есть одна нулевая мода (константная функция), соответствующая остаточной глобальной калибровке; поэтому в детерминантах нужно использовать $\det'$ и одновременно делить на объём остаточной группы калибровки;
  \item в 1-форменном секторе есть гармоническая 1-форма вдоль $S^1$ (пропорциональная $d\theta$), отвечающая плоской связности (Wilson line); она не должна входить в бесконечномерный гауссов интеграл, а учитывается конечномерным интегралом по модулю плоских связностей.
\end{itemize}
В $U(1)$-теории большие калибровки на $S^1$ отождествляют $a\sim a+n$, $n\in\mathbb{Z}$, если
\begin{equation}
  A_{\mathrm{harm}}=a\,d\theta.
\end{equation}
Следовательно, интеграл по модулю берётся по фундаментальной области $a\in[0,1)$ (или эквивалентно по одному классу голономии), и даёт конечный, $\mu$-независимый множитель.

\paragraph{Spin-структуры и периодичности (фиксация как часть постановки).}
Так как $\mathbb{RP}^3$ допускает две spin-структуры, спектральные объекты, содержащие оператор Дирака, зависят от выбора этой структуры.
Чтобы исключить эту свободу как ``ручку'', в постановке фиксируется продолжаемая (bounding) spin-структура на $\mathbb{RP}^3$.
Окружность $S^1$ трактуется как пространственная KK-окружность, поэтому фермионные поля берутся периодическими по $S^1$ (в отличие от термального KMS-круга, где естественна антипериодичность).

\subsection{ζ-регуляризация и масштаб $\mu$}
Для положительного самосопряжённого эллиптического оператора $\mathcal{O}$ с дискретным спектром определяем
\begin{equation}
  \zeta_{\mathcal{O}}(s)=\sum_{n}\frac{d_n}{\lambda_n^s},
  \qquad
  \ln\det\mathcal{O}:=-\zeta'_{\mathcal{O}}(0).
\end{equation}
При явном введении масштаба $\mu$ удобно фиксировать конвенцию
\begin{equation}
  \ln\det\!\left(\frac{\mathcal{O}}{\mu^2}\right)=-\zeta'_{\mathcal{O}}(0)-\zeta_{\mathcal{O}}(0)\,\ln\mu^2.
\end{equation}
На 4-мерном компактном фоне (для операторов лапласовского типа) ζ-регуляризация связана с коэффициентом теплового ядра $a_4$ формулой
\begin{equation}
  \zeta_{\mathcal{O}}(0)=a_4(\mathcal{O})-\dim\ker(\mathcal{O}).
\end{equation}
Это подчёркивает, что корректное обращение с $\det'$ (то есть с нулевыми модами) является частью определения схемы и критично для контроля $\mu$-зависимости.

\section{Сводный вывод кандидатной формулы для $\alpha^{-1}$ (самодостаточная версия)}
Цель этого приложения --- собрать в одном месте \emph{минимальную} цепочку утверждений, которая делает текст самодостаточным без внешних файлов.

\subsection{Геометрическое (локальное) ядро $S_{\mathrm{loc}}$}
Далее в этом приложении под «локальным ядром» понимается безразмерная константа
(в планковских единицах и при фиксированном $R=1$), которая фигурирует в численных sanity-check:
\begin{equation}
  S_{\mathrm{loc}}:=4\pi^3+\pi^2.
\end{equation}
Зависимость от радиуса $R$ (и связанный с ней размерный анализ) обсуждается отдельно через разложение вида \eqref{eq:alpha_candidate_appendix}; здесь же цель — зафиксировать именно то ядро, которое вычитается в ``budget check''.

\subsection{Постановочные леммы (что фиксируется как часть определения)}
\begin{itemize}
  \item \textbf{Фермионный сектор и накрытие $S^3$.}
  При фиксированной продолжаемой (bounding) spin-структуре на $\mathbb{RP}^3=S^3/\mathbb{Z}_2$ спиноры удобно описываются как $\mathbb{Z}_2$-эквивариантные спиноры на $S^3$; на уровне ζ-функций это реализуется проектором на нужный сектор, что позволяет выражать детерминанты через спектральные данные на накрытии. В рамках настоящей работы это используется как механизм, приводящий к объёмному множителю, пропорциональному $\mathrm{Vol}(S^3\times S^1)$.\cite{bar,ikeda}

  \item \textbf{Нормировка $U(1)$ и фактор $\mathrm{Vol}(S^1)$.}
  Бозонный множитель $\mathrm{Vol}(S^1)$ поглощается в определении 4D-константы связи после KK-редукции (эквивалентно: фиксируется комбинация $g_5^2/\mathrm{Vol}(S^1)$ нормировкой минимального заряда и периодичностью голономии по $S^1$).
\end{itemize}

\subsection{Систола $\pi R$ как ИК-масштаб и два сектора $\mathrm{Hol}=\pm 1$}
Так как $\pi_1(\mathbb{RP}^3)=\mathbb{Z}_2$, плоские $U(1)$-связности на $\mathbb{RP}^3$ имеют два дискретных класса голономии $\mathrm{Hol}=\pm 1$.
Геометрический масштаб, отвечающий генератору $\pi_1$, задаётся систолой $L_{\mathrm{sys}}(\mathbb{RP}^3)=\pi R$.
В данной версии текста $L_{\mathrm{sys}}$ рассматривается не как линейное слагаемое в $\alpha^{-1}$, а как естественный ИК-масштаб (граница применимости чисто локального разложения/сепарации секторов).

\paragraph{Что именно нужно вывести (в терминах $Z$).}
Рассмотрим абелеву теорию на $M_4\times K$ и зафиксируем калибровку (как в предыдущем приложении).
Для внутреннего фактора $K=\mathbb{RP}^3\times S^1$ существуют два дискретных плоских сектора по компоненте на $\mathbb{RP}^3$ (голономия $\pm 1$ вдоль генератора $\pi_1$).
Функциональный интеграл (после калибровочной фиксации) формально раскладывается в сумму по этим секторам:
\begin{equation}
  Z
  =
  \sum_{\mathrm{Hol}=\pm 1} Z_{\mathrm{Hol}},
  \qquad
  Z_{\mathrm{Hol}}:=\int_{A\in\mathcal{A}_{\mathrm{Hol}}/\mathcal{G}}\!\mathcal{D}A\,\exp\bigl(-S[A]\bigr)\,\times(\text{ghost}).
\end{equation}
Тогда вклад нетривиального сектора естественно параметризовать через \emph{безразмерную} разность эффективных действий:
\begin{equation}\label{eq:DeltaGamma_top_def}
  \Delta\Gamma_{\mathrm{top}}
  :=
  -\ln\frac{Z_{\mathrm{Hol}=-1}}{Z_{\mathrm{Hol}=+1}}.
\end{equation}
В этой работе $\delta_{\mathrm{top}}$ трактуется как секторная безразмерная величина
$\delta_{\mathrm{top}}:=-\ln(Z_{-}/Z_{+})$, которую требуется вычислить в выбранной постановке.

\emph{Замечание (об альтернативной физической гипотезе).}
Один из возможных сценариев, совместимых с размерным анализом 1-loop, состоит в том, что вклад нетривиального сектора в 4D-нормировку оказывается экспоненциально подавленным при больших $R$,
\begin{equation}
  \delta_{\mathrm{top}}(R)\sim \exp\bigl(-L_{\mathrm{sys}}/\ell_{\ast}\bigr)=\exp\bigl(-\pi R/\ell_{\ast}\bigr)
\end{equation}
для некоторого микроскопического масштаба $\ell_{\ast}$.
Однако численные ``budget check'' в текущей версии проекта используют другую рабочую гипотезу: что ведущий секторный масштаб может быть порядка $\pi$ (а малая разность $\mathrm{Rem}-\pi$ должна воспроизводиться конечной 1-loop частью и/или уточнением секторного отношения $Z_{-}/Z_{+}$).

\paragraph{Почему классическое действие не различает сектора.}
Для плоских связностей $F=0$, поэтому чисто Maxwell-овская часть действия равна нулю:
\begin{equation}
  S_{\mathrm{Maxwell}}[A_{\mathrm{flat}}]=0.
\end{equation}
Следовательно, если различие между секторами существует, оно должно возникать из:
(i) меры/якобиана при разложении поля и калибровочной фиксации,
(ii) 1-loop детерминантов операторов (Maxwell+ghost, и при наличии --- Dirac), скрученных на соответствующее плоское $U(1)$-расслоение.

\paragraph{Редукция к отношению скрученных ζ-детерминантов (план вычисления).}
В приближении 1-loop естественно ожидать структуру вида
\begin{equation}\label{eq:DeltaGamma_top_det_scheme}
  \Delta\Gamma_{\mathrm{top}}
  =
  \frac12\Bigl[\ln\det'\Delta^{(-)}_1-\ln\det'\Delta^{(+)}_1\Bigr]
  -\Bigl[\ln\det'\Delta^{(-)}_0-\ln\det'\Delta^{(+)}_0\Bigr]
  \;(+\;\text{при наличии }\Psi\text{ ещё }-\bigl[\ln\det D^{(-)}-\ln\det D^{(+)}\bigr]),
\end{equation}
где верхний индекс $\pm$ означает скручивание на плоский сектор $\mathrm{Hol}=\pm 1$.
Используя ζ-регуляризацию,
\begin{equation}
  \ln\det \mathcal{O} =-\zeta'_{\mathcal{O}}(0),
\end{equation}
задача сводится к вычислению разностей $\zeta'_{\mathcal{O}^{(-)}}(0)-\zeta'_{\mathcal{O}^{(+)}}(0)$ для скрученных операторов на $K$.

\paragraph{Где может появиться число $\pi R$ (и почему это не автоматическая истина).}
Топологическая структура задаёт два дискретных класса, а геометрия даёт естественную длину цикла $L_{\mathrm{sys}}=\pi R$.
Однако из одной лишь дискретности $\mathrm{Hol}=\pm 1$ не следует, что $\Delta\Gamma_{\mathrm{top}}$ обязан быть равен $\pi R$:
как показывают стандартные топологические инварианты (η-инвариант, Ray--Singer torsion), типичные ответы выражаются через $\log 2$ или дают ноль.
Поэтому в тексте ``$\pi R$'' следует понимать как \emph{обозначение цели расчёта}: нужно найти точную (и схемно устойчивую) величину \eqref{eq:DeltaGamma_top_def} и проверить, совпадает ли она с $\pi R$ в выбранной постановке.

\paragraph{Промежуточные строгие факты (и почему они не дают $\pi$ автоматически).}
Чтобы избежать скрытых допущений, фиксируем два стандартных источника ``топологических констант'' в 1-loop эффективных действиях на 3-многообразиях:

\begin{itemize}
  \item \textbf{η-инвариант.}
  Фазовая часть детерминанта фермионного оператора на нечётномерном многообразии может быть выражена через η-инвариант (APS-структура).
  Однако для $\mathbb{RP}^3$ в базовой (нескрученной) постановке спектр Дирака симметричен, и η-инвариант зануляется; следовательно, ``чистая'' η-механика не генерирует вклад масштаба $\pi R$.

  \item \textbf{Ray--Singer torsion.}
  В абелевых теориях и/или топологических редукциях отношение детерминантов на формах часто выражается через аналитическую (комбинаторную) торсию.
  Для нетривиального $U(1)$-представления $\pi_1(\mathbb{RP}^3)=\mathbb{Z}_2$ естественным образом возникает торсионный множитель порядка $\log 2$, а не $\pi$.
\end{itemize}

\paragraph{Численные проверки (на уровне toy-моделей).}
Для контроля того, что именно можно ожидать от ``универсальных'' топологических/спектральных механизмов, мы выполнили два простых воспроизводимых теста.

\begin{itemize}
  \item \textbf{Торсионный масштаб для нетривиального представления $\rho(-1)=-1$.}
  Скрипт \texttt{scripts/pi\_topology\_check.py} (Python 3.10, Linux) даёт
  \begin{equation}
    \log|1-\rho(-1)|=\log 2\approx 0.6931471805599453,
  \end{equation}
  что подчёркивает: естественный ``строгий'' топологический масштаб в абелевой постановке --- порядка $\log 2$, а не $\pi$.

  \item \textbf{Скрученный heat-kernel для скалярного лапласиана на $S^3$ (модель $\mathbb{RP}^3=S^3/\mathbb{Z}_2$).}
  Скрипт \texttt{scripts/pi\_sector\_heatkernel\_scalar\_S3.py} вычисляет (в toy-модели)
  \begin{equation}
    \Delta\log\det\;\sim\;-
    \int d\log t\,\sum_{n\ge 0}(-1)^n d_n e^{-t\lambda_n}
    \;\approx\;-4.30452474678
    \qquad (R=1,\;t\in[10^{-3},50],\;n\le 8000),
  \end{equation}
  причём результат устойчив по $n_{\max}$ в указанном диапазоне.
  Это демонстрирует, что ``голая'' $\mathbb{Z}_2$-скрутка в простейших спектральных объектах даёт конечную константу, но не выделяет автоматически значение $\pi$.

  \item \textbf{Toy Maxwell+ghost на $S^3/\mathbb{Z}_2$ (скрученные heat-trace для 0- и 1-форм).}
  Скрипт \texttt{scripts/pi\_sector\_heatkernel\_maxwell\_ghost\_S3.py} реализует toy-версию калибровочно-инвариантной комбинации (Maxwell+ghost) на модельной геометрии $\mathbb{RP}^3=S^3/\mathbb{Z}_2$ через скрученные следы вида $\mathrm{Tr}(g\,e^{-t\Delta_p})$.
  В этой постановке получено
  \begin{equation}
    I_0\approx -4.30452474678,
    \qquad
    I_1\approx -6.38114822702,
    \qquad
    \Delta\Gamma_{\mathrm{toy}}:=\tfrac12 I_1-I_0\approx 1.11395063327,
  \end{equation}
  при параметрах $R=1$, $t\in[10^{-3},50]$ и $n\le 8000$, причём результат устойчив по $n_{\max}$.
  Таким образом, даже в toy-версии, наиболее близкой по структуре к Maxwell+ghost, ``универсальная'' $\mathbb{Z}_2$-скрутка приводит к константе порядка 1, а не автоматически к $\pi$.

  \item \textbf{Проверка отсутствия линейного тренда по $R$ (scale-invariant $\tau$-окно).}
  Чтобы отделить реальную зависимость от масштаба от артефакта численного усечения, важно сравнивать ответы при одном и том же безразмерном параметре $\tau:=t/R^2$.
  Скрипт \texttt{scripts/pi\_sector\_R\_scan\_maxwell\_ghost\_S3.py} выполняет два прогона:
  (A) фиксирует окно по $t$ и получает сильную зависимость $\Delta\Gamma_{\mathrm{toy}}(R)$, и
  (B) фиксирует окно по $\tau$ (то есть берёт $t\in[R^2\tau_{\min},R^2\tau_{\max}]$) и обнаруживает
  \begin{equation}
    \Delta\Gamma_{\mathrm{toy}}(R)\approx 1.11395063327
  \end{equation}
  с наклоном линейной регрессии $b\sim 10^{-13}$ на диапазоне $R\in[0.7,1.3]$.
  Тем самым ``линейность по $R$'' в режиме (A) является чистым следствием фиксированного $t$-окна, а в масштабно-эквивалентной постановке (B) toy-ответ становится $R$-независимым (как и ожидается по размерному анализу 1-loop объектов).

  \item \textbf{Численная ``смета'' для $\alpha^{-1}$ (budget check после отделения локального ядра).}
  Если выделить локальную (объёмную) часть геометрического ядра как
  \begin{equation}
    S_{\mathrm{loc}}:=4\pi^3+\pi^2, 
  \end{equation}
  то остаток, который должен быть объяснён суммой топологического сектора и конечной 1-loop части, равен
  \begin{equation}
    \mathrm{Rem}:=\alpha^{-1}_{\mathrm{exp}}-S_{\mathrm{loc}}\approx 3.141288054711.
  \end{equation}
  Численно он близок к $\pi$, но отличается на величину
  \begin{equation}
    \mathrm{Rem}-\pi\approx -3.045988783974\times 10^{-4}.
  \end{equation}
  Воспроизводимые скрипты показывают, что эта разность по порядку величины согласуется с масштабом $-1/(24S)$ при $S\approx S_{\mathrm{loc}}+\pi$ (см. \texttt{scripts/alpha\_budget\_check.py}, \texttt{scripts/alpha\_budget\_candidates.py}, \texttt{scripts/alpha\_recompose\_scan.py}).
\end{itemize}

Эти пункты важны методологически: если целевой вклад $\Delta\Gamma_{\mathrm{top}}$ действительно пропорционален $\pi R$, то он должен возникать не из ``универсальных'' 3D-инвариантов (η/торсия) и не из простейших toy-скруток, а из конкретной 4D-постановки и её калибровочно-инвариантной комбинации скрученных операторов на $K=\mathbb{RP}^3\times S^1$ (включая правило $\det'$ и обращение с гармоническими модами).

\paragraph{Промежуточный результат: что даёт ``чистая'' Maxwell+ghost торсия.}
В 3-мерном контексте (и для компактного $M^3$ без границы) калибровочно-инвариантная 1-loop комбинация детерминантов Maxwell+ghost может быть связана с аналитической торсией Рея--Зингера $T_{\mathrm{RS}}(M^3,\rho)$ для соответствующего унитарного представления $\rho$ (скрученного плоского расслоения).
В силу теоремы Cheeger--M{"u}ller, $T_{\mathrm{RS}}$ совпадает с комбинаторной торсией Рейдемайстера и потому даёт \emph{действительно} топологический инвариант.\cite{cheeger,muller}
Для $\pi_1(\mathbb{RP}^3)=\mathbb{Z}_2$ и нетривиального унитарного представления $\rho(-1)=-1$ такая торсия естественным образом приводит к масштабу $\log 2$.
Тем самым, если целевой вклад $\Delta\Gamma_{\mathrm{top}}$ равен $\pi R$, то он не может быть объяснён только через ``чистую'' 3D-торсию Maxwell+ghost: необходимо, чтобы (i) играла роль 4D-ренормировка, (ii) учитывались дополнительные сектора (например, фермионный детерминант/мера) и/или (iii) вклад возникал из конкретной 4D-постановки на $\mathbb{RP}^3\times S^1$, а не из универсального 3D-инварианта.

\paragraph{Перспектива: TQFT/локализация.}
Один из путей строгого вывода --- связать \eqref{eq:DeltaGamma_top_def} с топологической теорией на 3-многообразии (например, с $U(1)$ Черна--Саймонса на $\mathbb{RP}^3$) и локализовать вклад на плоских связностях; далее нужно показать, как именно этот объект встраивается в 4D-ренормусловие для $\alpha(0)$.\cite{witten_cs}

\subsection{Касимировский (ζ-регуляризованный) остаток и коэффициент $1/24$}
На одномерном круге стандартная ζ-регуляризация даёт $\zeta_R(-1)=-1/12$, т.е.
\begin{equation}
  \sum_{n\ge 1} n
  \;\overset{\zeta}{=}\;
  -\frac{1}{12}.
\end{equation}
В калибровочно-инвариантной комбинации Maxwell+ghost возникает «половинный» остаток, что даёт кандидата $\kappa_{\mathrm{Cas}}=1/24$.
Строгость требует контроля схемозависимости через heat-kernel/$a_4$ и $\zeta_{\mathcal{O}}(0)$ для соответствующей комбинации операторов.\cite{gilkey}

\subsection{Сопоставление с $\alpha(0)$ и RG-мост}
В рамках фиксированной постановки корректная по размерностям цель формулируется как разложение
\begin{equation}\label{eq:alpha_candidate_appendix}
  \alpha^{-1}(0)
  =
  \frac{\mathrm{Vol}(K(R))}{g_5^2}
  +\Delta^{\mathrm{finite}}_{\mathrm{1\text{-}loop}}(K(R))
  +\delta_{\mathrm{top}}(K(R)),
\end{equation}
где $\Delta^{\mathrm{finite}}_{\mathrm{1\text{-}loop}}$ --- конечная часть калибровочно-инвариантной 1-loop комбинации детерминантов (Maxwell+ghost, с фиксированным правилом $\det'$ и после вычитания локальных контрчленов), а $\delta_{\mathrm{top}}$ параметризует вклад нетривиального сектора голономии ($\mathrm{Hol}=-1$) и связана с безразмерной величиной $-\ln(Z_{-}/Z_{+})$.
Систола $L_{\mathrm{sys}}=\pi R$ задаёт естественный ИК-масштаб для разделения секторов, но не обязана порождать линейный 1-loop вклад.

Интерпретация делается как условие ИК-нормировки $\alpha(0)$; перенос на другие масштабы $\mu$ выполняется стандартным ренормгрупповым бегом КЭД/СМ.
Минимальная проверка: при $\mu>m_e$ выполняется
\begin{equation}
  \mu\frac{d\alpha}{d\mu}=\frac{2}{3\pi}\,\alpha^2+\cdots, 
\end{equation}
что эквивалентно
\begin{equation}
  \alpha^{-1}(\mu)=\alpha^{-1}(0)-\frac{2}{3\pi}\ln\frac{\mu}{m_e}+\cdots.
\end{equation}

\begin{thebibliography}{12}
\bibitem{gilkey}
P.~B.~Gilkey.
\textit{Invariance Theory, the Heat Equation, and the Atiyah--Singer Index Theorem}.
CRC Press, 2nd ed., 1995.

\bibitem{bar}
C.~B{"a}r.
The Dirac operator on space forms of positive curvature.
\textit{J. Math. Soc. Japan} \textbf{48} (1996), 69--83.

\bibitem{ikeda}
A.~Ikeda.
Spectra and eigenforms of the Laplacian on $S^n$ and $P^n(\mathbb{C})$.
\textit{Osaka J. Math.} \textbf{15} (1978), 515--546.

\bibitem{birrell_davies}
N.~D.~Birrell, P.~C.~W.~Davies.
\textit{Quantum Fields in Curved Space}.
Cambridge University Press, 1982.

\bibitem{peskin}
M.~E.~Peskin, D.~V.~Schroeder.
\textit{An Introduction to Quantum Field Theory}.
Westview Press, 1995.

\bibitem{witten_cs}
E.~Witten.
Quantum field theory and the Jones polynomial.
\textit{Commun. Math. Phys.} \textbf{121} (1989), 351--399.

\bibitem{cheeger}
J.~Cheeger.
Analytic torsion and the heat equation.
\textit{Ann. of Math.} \textbf{109} (1979), 259--322.

\bibitem{muller}
W.~M{"u}ller.
Analytic torsion and R-torsion of Riemannian manifolds.
\textit{Adv. Math.} \textbf{28} (1978), 233--305.
\end{thebibliography}

\end{document}